    % Template KLTN cho SV trường ĐHKHTN
% Liên hệ: nqminh@fit.hcmus.edu.vn
% Last update: 08/06/2016

% Chú ý: đọc các phần chú ý đóng khung của file này và chỉnh lại cho phù hợp.
% Trước khi build, xóa hết các file được tạo ra trong quá trình build trước đó, và build theo thứ tự: BIB > PDF > PDF.
% Nếu cập nhật tài liệu tham khảo, cũng cần build lại theo cách trên.

\documentclass[oneside,a4paper,14pt]{extreport}

% Font tiếng Việt
\usepackage[T5]{fontenc}
\usepackage[utf8]{inputenc}
\usepackage{amsmath}
\usepackage{amssymb}
\DeclareTextSymbolDefault{\DH}{T1}

% Tài liệu tham khảo
\usepackage[
	sorting=nty,
	backend=bibtex,
	defernumbers=true]{biblatex}
\usepackage[unicode]{hyperref} % Bookmark tiếng Việt
\addbibresource{References/references.bib}

\makeatletter
\def\blx@maxline{77}
\makeatother

% Chèn hình, các hình trong luận văn được để trong thư mục Images/
\usepackage{graphicx}
\graphicspath{ {Images/} }

% Chèn và định dạng mã nguồn
\usepackage{listings}
\usepackage{color}
\definecolor{codegreen}{rgb}{0,0.6,0}
\definecolor{codegray}{rgb}{0.5,0.5,0.5}
\definecolor{codepurple}{rgb}{0.58,0,0.82}
\definecolor{backcolour}{rgb}{0.95,0.95,0.92}
\lstdefinestyle{mystyle}{
    backgroundcolor=\color{backcolour},   
    commentstyle=\color{codegreen},
    keywordstyle=\color{magenta},
    numberstyle=\tiny\color{codegray},
    stringstyle=\color{codepurple},
    basicstyle=\footnotesize,
    breakatwhitespace=false,         
    breaklines=true,                 
    captionpos=b,                    
    keepspaces=true,                 
    numbers=left,                    
    numbersep=5pt,                  
    showspaces=false,                
    showstringspaces=false,
    showtabs=false,                  
    tabsize=2
}
\lstset{style=mystyle}

% Chèn và định dạng mã giả
\usepackage{amsmath}
\usepackage{algorithm}
\usepackage[noend]{algpseudocode}
\makeatletter
\def\BState{\State\hskip-\ALG@thistlm}
\makeatother

% Bảng biểu
\usepackage{multirow}
\usepackage{array}
\newcolumntype{L}[1]{>{\raggedright\let\newline\\\arraybackslash\hspace{0pt}}m{#1}}
\newcolumntype{C}[1]{>{\centering\let\newline\\\arraybackslash\hspace{0pt}}m{#1}}
\newcolumntype{R}[1]{>{\raggedleft\let\newline\\\arraybackslash\hspace{0pt}}m{#1}}

% Đổi tên mặc định
\renewcommand{\chaptername}{Chương}
\renewcommand{\figurename}{Hình}
\renewcommand{\tablename}{Bảng}
\renewcommand{\contentsname}{Mục lục}
\renewcommand{\listfigurename}{Danh sách hình}
\renewcommand{\listtablename}{Danh sách bảng}
\renewcommand{\appendixname}{Phụ lục}

% Định nghĩa
\newtheorem{definition}{Định nghĩa}
% Vẽ cây
\usepackage{tikz}
\usetikzlibrary{calc}
\usetikzlibrary{trees}

% Dãn dòng 1.5
\usepackage{setspace}
\onehalfspacing

% Thụt vào đầu dòng
\usepackage{indentfirst}

% Canh lề
\usepackage[
  top=30mm,
  bottom=25mm,
  left=30mm,
  right=20mm,
  includefoot]{geometry}
  
% Trang bìa
\usepackage{tikz}
\usetikzlibrary{calc}
\newcommand\HRule{\rule{\textwidth}{1pt}}
% ========================================================================================= %
% PACKAGE CUSTOMIZE
% pdf
\usepackage{pdfpages}

\usepackage[strings]{underscore}
\usepackage{url}

% Math
\usepackage{amsfonts}

\usepackage{stmaryrd}
\newcommand*\concat{\mathbin{\|}} % parallel

% Vẽ cây
\usetikzlibrary{trees}

% subfigure
\usepackage{subcaption}
\usepackage[justification=centering]{caption}

% Bảng thuật ngữ
\usepackage[acronym,nomain]{glossaries}
% ========================================================================================= %
% CHÚ Ý: Thông tin chung về KLTN - sinh viên điền vào đây để tự động update các trang khác  %
% ========================================================================================= %
\newcommand{\tenSV}{Phan~Minh~Tâm~-~Hoàng~Minh~Thanh} % Dấu ~ là khoảng trắng không được tách (các chữ nối với nhau bằng dấu ~ sẽ nằm cùng 1 dòng
\newcommand{\mssv}{18424059-18424062}
\newcommand{\tenKL}{ĐỰ~ĐOÁN~LIÊN~KẾT~TRONG~ĐỒ~THỊ~TRI~THỨC} % Chú ý dấu ~ trong tên khóa luận
\newcommand{\tenGVHD}{ThS.~Lê~Ngọc~Thành}
\newcommand{\tenBM}{BM. Khoa~Học~Máy~Tính}

\begin{document}

\begin{titlepage}

\begin{center}
%ĐẠI HỌC QUỐC GIA THÀNH PHỐ HỒ CHÍ MINH\\
TRƯỜNG ĐẠI HỌC KHOA HỌC TỰ NHIÊN\\
\textbf{KHOA CÔNG NGHỆ THÔNG TIN}\\[2cm]


{ \Large \bfseries Hoàng Minh Thanh - Phan Minh Tâm\\[2cm] } 

%Tên đề tài Khóa luận tốt nghiệp/Đồ án tốt nghiệp

{ \Large \bfseries DỰ ĐOÁN LIÊN KẾT TRONG\\ĐỒ THỊ PHỨC\\KNOWLEDGE GRAPH EMBEDDING FOR LINKING PREDICTION\\[2cm] } 


%Chọn trong các dòng sau
\large KHÓA LUẬN TỐT NGHIỆP CỬ NHÂN\\
%\large ĐỒ ÁN TỐT NGHIỆP CỬ NHÂN\\
%\large THỰC TẬP TỐT NGHIỆP CỬ NHÂN\\
%Đưa vào dòng này nếu thuộc chương trình Chất lượng cao, hoặc lớp Cử nhân tài năng
\large CHƯƠNG TRÌNH HOÀN CHỈNH ĐẠI HỌC\\
%\large CHƯƠNG TRÌNH CHẤT LƯỢNG CAO\\
%\large CHƯƠNG TRÌNH CỬ NHÂN TÀI NĂNG\\[2cm]

\begin{tikzpicture}[remember picture, overlay]
  \draw[line width = 2pt] ($(current page.north west) + (2cm,-2cm)$) rectangle ($(current page.south east) + (-1.5cm,2cm)$);
\end{tikzpicture}

\vfill
Tp. Hồ Chí Minh, tháng 09/2020

\end{center}

\pagebreak



\begin{center}

TRƯỜNG ĐẠI HỌC KHOA HỌC TỰ NHIÊN\\
\textbf{KHOA CÔNG NGHỆ THÔNG TIN}\\[2cm]


{\large \bfseries Hoàng Minh Thanh - 18424062\\} 
{\large \bfseries Phan Minh Tâm - 18424059\\[2cm]}

%Tên đề tài Khóa luận tốt nghiệp/Đồ án tốt nghiệp

{ \Large \bfseries DỰ ĐOÁN LIÊN KẾT TRONG\\ĐỒ THỊ PHỨC\\KNOWLEDGE GRAPH EMBEDDING FOR LINKING PREDICTION\\[2cm] } 


%Chọn trong các dòng sau
\large KHÓA LUẬN TỐT NGHIỆP CỬ NHÂN\\
%\large ĐỒ ÁN TỐT NGHIỆP CỬ NHÂN\\
%Đưa vào dòng này nếu thuộc chương trình Chất lượng cao, hoặc lớp Cử nhân tài năng
\large CHƯƠNG TRÌNH HOÀN CHỈNH ĐẠI HỌC\\[2cm]
%\large CHƯƠNG TRÌNH CHẤT LƯỢNG CAO\\[2cm]
%\large CHƯƠNG TRÌNH CỬ NHÂN TÀI NĂNG\\[2cm]

\textbf{GIÁO VIÊN HƯỚNG DẪN}\\
Th.S Lê Ngọc Thành\\

\begin{tikzpicture}[remember picture, overlay]
  \draw[line width = 2pt] ($(current page.north west) + (2cm,-2cm)$) rectangle ($(current page.south east) + (-1.5cm,2cm)$);
\end{tikzpicture}

\vfill
Tp. Hồ Chí Minh, tháng 09/2020

\end{center}

\end{titlepage}
% Sasu trang Title, các bạn chèn nhận xét gủa GVHD và GVPB. Nhận xét sẽ được giáo vụ phát sau buổi bảo vệ để các bạn đóng quyển.

\pagenumbering{roman} % Đánh số i, ii, iii, ...

\addcontentsline{toc}{chapter}{Lời cam đoan}
\chapter*{Lời cam đoan}
\label{reassurances}

Tôi xin cam đoan đây là công trình nghiên cứu của riêng chúng tôi. Các số liệu và kết quả nghiên cứu trong luận văn này là trung thực và không trùng lặp với các đề tài khác.

\addcontentsline{toc}{chapter}{Lời cảm ơn}
\chapter*{Lời cảm ơn}
\label{thanks}

Tôi xin chân thành cảm ơn ...

\addcontentsline{toc}{chapter}{Đề cương chi tiết}
%%Đây là template dùng cho đề cương đề tài tốt nghiệp
%Khoa Công nghệ Thông tin
%Trường Đại học Khoa học Tự nhiên, ĐHQG-HCM

%Liên hệ về mẫu LaTEX này: Thầy Bùi Huy Thông (bhthong@fit.hcmus.edu.vn)

\documentclass{article}[14pt]
\usepackage[utf8]{vietnam}
\usepackage{enumerate}
\usepackage{enumitem}
\usepackage{multicol}
\usepackage{listings}
\usepackage[left=2cm,right=2cm,top=2.5cm,bottom=2.5cm]{geometry}
\usepackage{verbatim}
\usepackage{graphicx}
\usepackage{url}
\usepackage{fancyhdr}
\usepackage{fancybox,framed}
\linespread{1.2}
\usepackage{lastpage}
\usepackage{floatrow}
\usepackage{floatrow}
\pagenumbering{arabic}
\usepackage{dirtytalk}
%\pagestyle{fancy}
\newfloatcommand{capbtabbox}{table}[][\FBwidth]

\usepackage{blindtext}
\usepackage{titlesec}
\usepackage[nottoc]{tocbibind}

\titleformat*{\section}{\LARGE\bfseries}
\titleformat*{\subsection}{\Large\bfseries}
\titleformat*{\subsubsection}{\large\bfseries}
%\addbibresource{ref.bib}


\begin{document}
    \begin{figure}[h]
        \begin{floatrow}
        \ffigbox{\includegraphics[scale = 0.1]{logo.png}}  
        {%
    
        }
        \capbtabbox{
            \begin{tabular}{l}
            \multicolumn{1}{c}{\textbf{\begin{tabular}[c]{@{}c@{}}TRƯỜNG ĐẠI HỌC KHOA HỌC TỰ NHIÊN\\KHOA CÔNG NGHỆ THÔNG TIN\end{tabular}}} \\ \\ \\
            \end{tabular}
        }
        {%
    
        }
        \end{floatrow}
    \end{figure}
    
    \begin{center}
        
        %Xác định loại đề tài tốt nghiệp tương ứng: Khóa luận, Thực tập, Đồ án
        \textbf{\Large ĐỀ CƯƠNG KHOÁ LUẬN TỐT NGHIỆP} \\ 
    \end{center}
    
    %\vspace{.5cm}
    
    \begin{center}
    %Tên đề tài phải VIẾT HOA
        
        \textbf{\huge Dự Đoán Liên Kết Trong Đồ Thị Tri Thức} 
        \\
        
    %Tên đề tài bằng tiếng Anh (nếu có)
    \vspace{.5cm}
        \textit{\textbf{\Large (Knowledge Graph Embedding for Link Prediction)}}
    \end{center}
    
    \vspace{.5cm}
    
    \Large
    \section{THÔNG TIN CHUNG}
    \begin{itemize}[label = {}]
        
        \item \textbf{Người hướng dẫn:} 
        %Thể hiện dạng: <Chức danh> <Họ và tên> (<Đơn vị công tác>)
        \begin{itemize}
            \item Ths. Lê Ngọc Thành (Khoa Công nghệ Thông tin)
        \end{itemize}{}
    
        
        \item \textbf{Nhóm Sinh viên thực hiện:}
        
        %Thể hiện dạng: <Họ và tên sinh viên> (MSSV: )
        \begin{enumerate}
        
            \item Phan Minh Tâm (18424059)
            \item Hoàng Minh Thanh (18424062)
        \end{enumerate}

       %Chọn loại thích hợp
        \item \textbf{Loại đề tài:} Nghiên cứu
        
        \item \textbf{Thời gian thực hiện:} Từ \textit{04/2020} đến \textit{10/2020}
        
        
    \end{itemize}
    
    \section{NỘI DUNG THỰC HIỆN}
    {

    %Mỗi mục dưới đây phải viết ít nhất là 5 câu mô tả/giới thiệu.
    
    \subsection{Giới thiệu về đề tài}
    
    Đồ thị tri thức (Knowledge Graphs-KG) là các biểu diễn cấu trúc của thông tin thế giới thực. Do khả năng mô hình hóa dữ liệu có cấu trúc, phức tạp theo cách máy tính có thể dễ dàng \say{hiểu được}, KG hiện đang được sử dụng rộng rãi trong nhiều lĩnh vực khác nhau, từ trả lời câu hỏi đến truy xuất thông tin và các hệ thống có thể suy luận dựa trên nội dung đã có.
    Việc phát triển một KG có thể được thực hiện bằng cách trích xuất các sự kiện mới từ các nguồn bên ngoài hoặc bằng cách suy ra các sự kiện còn thiếu từ những sự kiện đã có trong KG. Phương pháp tiếp cận, được gọi là Dự đoán liên kết (Link Prediction-LP).
    
    \subsection{Mục tiêu đề tài}
    %Phần này mô tả mục tiêu thực hiện đề tài.
    Cùng với nhiều kỹ thuật trí tuệ nhân tạo phát triển mạnh gần đây, đề tài tập chung nghiên cứu vào các khía cạnh của bài toán LP trên KG như đặc trưng tập dữ liệu, thuật toán, thực nghiệm đánnh giá các phương pháp cũng như các kỹ thuật khác nhau cùng tìm hiểu xem liệu những đặc trưng gì của tập dữ liệu hoặc các thuật toán khác nhau ảnh hưởng tới khả năng khái quát hóa của mô hình.
    \subsection{Phạm vi của đề tài}
    
    LP là một lĩnh vực nghiên cứu ngày càng sôi nổi gần đây đã phát triển mạnh mẽ từ sự bùng nổ của các kỹ thuật trong trí tuệ nhân tạo như: máy học (machine learning) và kỹ thuật học sâu (deep learning). Đề tài sẽ tập trung nghiên cứu các mô hình LP sử dụng KG làm nền tảng để tìm hiểu các biểu diễn dữ liệu với số chiều thấp còn được gọi là Knowledge Graph Embeddings, sau đó sử dụng chúng để suy ra các sự kiện, quan hệ mới.
    
    \subsection{Cách tiếp cận dự kiến}
    
    %Có thể bổ sung hình ảnh vào để làm rõ phương pháp hoặc cách tiếp cận dự kiến.
    Với mong muốn được tìm hiểu và nghiên cứu các phương pháp trí tuệ nhân tạo được áp dụng vào LG từ cổ điển đến hiện đại đề tài sẽ nghiên cứu một phương pháp cổ điển dựa trên luật (rule base \cite{burl}) và phương pháp còn lại dựa vào phương pháp học sâu (learning attention \cite{gat}). Đối với phương pháp cổ điển thuật toán cố gắng đưa các luật Horn \cite{wiki:Hornclause} sau đó khái quát hóa chúng thành các luật tổng quát trên tập dữ liệu huấn luyện. Phương pháp còn lại dựa vào kỹ thuật nhúng \cite{BordesUGWY13} đồ thị tri thức vào một không gian số chiều thấp dựa trên các láng giềng lân cận \cite{gat} sau đó học trên không gian này để đưa ra các dự đoán về các cạnh của đồ thị.
    
    \subsection{Kết quả dự kiến của đề tài}
    Tìm hiểu các phương pháp khác thực nghiệm và so sánh kết quả với phương pháp đang nghiên cứu với các phương pháp khác trên bốn tập dữ liệu được nghiên cứu nhiều nhất là FB15k \cite{toutanova2015observed}, FB15k-273 \cite{toutanova2015observed}, WN188 \cite{BordesUGWY13}, WN188RR \cite{GEmbedding}
    
    \subsection{Kế hoạch thực hiện}
    \begin{itemize}
        \item 4/2020 - 6/2020: Tìm hiểu các kiến thức liên quan về mạng nơ-ron,  KG. Đọc các  tài liệu bài báo liên quan tới đề tài.
        \item 6/2020 - 8/2020: Hiện thực các phương pháp, thuât toán đã đề xuất và tìm hiểu.
        \item 8/2020 - 9/2020: Tìm hiểu các phương pháp khác thực nghiệm và so sánh kết quả với phương pháp nghiên cứu.
        \item 9/2020 - 10/2020: Báo cáo và bảo vệ khóa luận tốt nghiệp.
    \end{itemize}
    }
    %TÀI LIỆU TRÍCH DẪN
    %Đây là ví dụ
    \bibliographystyle{ieeetr}
    \bibliography{sample}
    \nocite{*}

    \begin{table}[h]
    \centering
        \begin{tabular}{p{7cm}p{7cm}}
        \textbf{\begin{tabular}[c]{@{}c@{}}\\XÁC NHẬN\\CỦA NGƯỜI HƯỚNG DẪN\\ \textit{(Ký và ghi rõ họ tên)}\end{tabular}} & \textbf{\begin{tabular}[c]{@{}c@{}}\textit{TP. Hồ Chí Minh, ngày 20 tháng 09 năm 2020}\\NHÓM SINH VIÊN THỰC HIỆN\\\textit{(Ký và ghi rõ họ tên}) \end{tabular}}
        \end{tabular}
    \end{table}
    
\end{document}              

% Mục lục, danh sách hình, danh sách bảng
\addcontentsline{toc}{chapter}{Mục lục}
\tableofcontents
\listoffigures
\listoftables

\addcontentsline{toc}{chapter}{Tóm tắt}
\include{Appendix/tomtat}

\clearpage

\pagenumbering{arabic} % Đánh số 1, 2, 3, ...

% Các chương nội dung
% Giới thiệu
%❖ Là chương mở rộng của tóm tắt
%❖ Nội dung bao gồm
%o Ngữ cảnh, vấn đề cần giải quyết là gì?
%o Vì sao vấn đề/bài toán quan trọng và thú vị?
%o Bài toán có gì khó? Vì sao cần phải giải quyết?
%o Những giải pháp (nghiên cứu, ứng dụng) đã giải quyết bài toán này?
%o Những giải pháp này có hạn chế, thiếu sót gì?
%o Giải pháp của bạn là gì? Kết quả thế nào?
%o Đóng góp của nghiên cứu/ứng dụng của bạn? 
%❖ Độ dài từ 3 đến 8 trang
%❖ Có thể chia ra các mục
%o Đặt vấn đề (Problem)
%o Mục tiêu (Objectives)
%o Giải pháp/Cách tiếp cận (An Approach)
%o Đóng góp (Contributions)
%o Bố cục (Outline)
%❖ Lưu ý, nội dung các chương sau phải khớp với mục tiêu, giải pháp, đóng góp ở chương này
\chapter{Giới thiệu}
\label{Chapter1}

Ngày nay đồ thị đã được ứng dụng vào mọi mặt của đời sống, với đồ thị về mạng xã hội (Facebook \cite{ugander2011anatomy}) thể hiện thông tin kết nối từng người với nhau, những nơi chúng ta đến, những thông tin chúng ta tương tác, đồ thị cũng được sử dụng làm cấu trúc trong hệ thống gợi ý video (Youtube \cite{baluja2008video}), trong mạng các chuyến bay, hệ thống định vị GPS, những tính toán khoa học hay thậm chí là kết nối não. Đồ thị tri thức của Google (Google's Knowledge Graph \cite{googlekg:2020}) được Google giới thiệu năm 2012 \cite{ji2020survey}, một loại đồ thị biểu diễn thông tin, là một trong những ứng dụng rõ ràng nhất về đồ thị tri thức cũng như cách dữ liệu được khai thác và biểu diễn trên đồ thị tri thức.

Khai thác đồ thị tri thức hiệu quả cung cấp cho người dùng hiểu sâu hơn về những gì đằng sau dữ liệu và từ đó có thể mang lại lợi ích cho nhiều ứng dụng trong thực tế. Tuy nhiên, trong thực tế, luôn có những tri thức mới được sinh ra mỗi ngày, thông tin thu được thường bị mất mát và không đầy đủ, từ đó nảy sinh ra vấn đề hoàn thiện đồ thị tri thức (knowledge graph completion) hay dự đoán liên kết (linking prediction) trong đồ thị tri thức.
Hầu hết các phương tiếp cận hiện nay là dự đoán một cạnh mới nối từ đỉnh này tới đỉnh khác. Với cách tiếp cận như vậy đồ thị tri thức có đầy đủ các tri thức nói cách khác làm cho đồ thị đày đặc nhờ tạo thêm cách cạnh nối. Nhưng như vậy chỉ mới giải quyết được vấn đề hoàn thành đồ thị, vấn đề thêm một (hoặc một lượng) tri thức mới vào đồ thị vẫn còn là một câu hỏi mở.

Hiện nay các bài toán liên quan đến hoàn thành đồ thị tri thức có hai cách tiếp cận chính là tối ưu hóa hàm mục tiêu tức là đưa ra dự đoán đựa trên có sai sót ít nhất như trong RuDiK\cite{ortona2018robust}, AMIE\cite{galarraga2015fast}, RuleN\cite{meilicke2018fine} liên quan tới các ứng dụng phân loại đỉnh, phân loại cạnh. Hoặc đưa ra một danh sách gồm \(k\) ứng viên với số điểm đại điện cho độ tin cậy giảm đần như trong các nghiên cứu TransE\cite{bordes2013translating}, ConvKB\cite{vu2019capsule} liên quan tới các hệ thống gợi ý (recommend system)...Cách tiếp cận của chúng tôi dựa trên cách tạo ra một danh sách gồm \(k\) ứng viên này.

Với mỗi cách tiếp cận trên có hai phương pháp chính đưa ra nghiên cứu một cách dựa trên luật như trong AnyBURL\cite{burl} hoặc dựa trên nhúng đồ thị như trong ConvE\cite{dettmers2017convolutional}, TransE\cite{bordes2013translating}, ComplEx \cite{trouillon2016complex}. Với mong muốn được tiếp cận các phương pháp có hệ thống nên chúng tôi chọn ở cả hai phương pháp để nghiên cứu thực hiện đề tài này. Đối với phương pháp dựa trên luật chúng tôi chọn phương pháp AnyBURL\cite{burl} còn với phương pháp dựa trên nhúng đồ thị chúng tôi chọn phương pháp KBAT \cite{nathani2019learning} sử dụng cơ chế chú  ý.

Đóng góp của chúng tôi trong phương pháp AnyBURL\cite{burl} gồm mã nguồn Python phương pháp AnyBURL. Cùng với đó chúng tôi cung cấp thêm hai chiến lược để thêm tri thức mới vào đồ thị mà chúng tôi gọi là online-to-offline là mở rộng của AnyBURL trong việc tạo ra các luật khi có một lượng (tập hợp) tri thức mới được thêm vào. Online-to-online sẽ tạo ngay các luật mới khi có một tri thức mới(cạnh) được thêm vào.

Với phương pháp dựa trên nhúng đồ thị chúng tôi sẽ trình bày lại về cơ chế chú ý (attention mechanisms \cite{vaswani2017attention}), cách cơ chế chú ý được áp dụng vào đồ thị tri thức bằng mô hình Mạng Đồ Thị Chú Ý (Graph Attention Network - GAT \cite{velivckovic2017graph}), mô hình KBAT \cite{nathani2019learning}, cũng như cải tiến của chúng tôi trên đồ thị dưa trên cải tiến mới nhất của cơ chế chú ý.
Đóng góp của chúng tôi trong phương pháp học sâu bao gồm mã nguồn trực tuyến trên Google Colab về cách cải tiến của mô hình chúng tôi, cũng như các phân tích của chúng tôi về các siêu tham số trong quá trình huấn luyện để đạt được kết quả tốt hơn.

% Hiện nay các bài toán liên quan đến hoàn thành đồ thị tri thức rất được quan tâm có bốn nhánh nghiên cứu chính như được nhắc đến trong nghiên cứu \cite{ampligraph} có hai nhánh nghiên cứu chính là nhúng đồ thị \cite{cai2018comprehensive} và dựa trên luật \cite{burl}. Với mong muốn được tiếp cận các phương pháp có hệ thống nên chúng tôi chọn ở mỗi nhánh một phương pháp làm chủ để chính cho các nghiên cứu và báo cáo này. Trong bài toán dự đự đoán liên kết cũng có hai phương pháp chính được đề xuất một là tối ưu hóa hàm mục tiêu đ


% Hiện nay các bài toán liên quan đến dự đoán liên kết đồ thị tri thức lớn rất được quan tâm có khoảng bốn nhánh nghiên cứu chính như được nhắc đến trong nghiên cứu \cite{ampligraph} một trong số đó là phương pháp dựa trên luật logic. Với mong muốn được tiếp cận các phương pháp từ đơn giản đến phức tạp nên chúng tôi chọn phương pháp này làm chủ để chính cho các báo cáo và nghiên cứu trong chương này. Với phương pháp này đưa ra một xếp hạng \(k\) ứng viên với một số điểm nhất định biểu thị cho sự chắc chán của dự đoán nó phù hợp với các hệ thống gợi ý(recommender system).

% Ngôn ngữ để viết và trình bày báo cáo khóa luận tốt nghiệp, đồ án tốt nghiệp, thực tập tốt nghiệp (sau đây gọi chung là báo cáo) là tiếng Việt hoặc tiếng Anh. 
% Trường hợp chọn ngôn ngữ tiếng Anh để viết và trình bày báo cáo,  sinh viên cần có đơn đề nghị, được cán bộ hướng dẫn (CBHD) đồng ý và nộp cho bộ phận Giáo vụ của Khoa vào thời điểm đăng ký đề tài để xin ý kiến.
% Báo cáo viết và trình bày bằng tiếng Anh phải có bản tóm tắt viết bằng tiếng Việt.


%Tóm tắt luận văn được trình bày nhiều nhất trong 24 trang in trên hai mặt giấy, cỡ chữ Times New Roman 11 của hệ soạn thảo Winword hoặc phần mềm soạn thảo Latex đối với các chuyên ngành thuộc ngành Toán.

%Mật độ chữ bình thường, không được nén hoặc kéo dãn khoảng cách giữa các chữ.
%Chế độ dãn dòng là Exactly 17pt.
%Lề trên, lề dưới, lề trái, lề phải đều là 1.5 cm.
%Các bảng biểu trình bày theo chiều ngang khổ giấy thì đầu bảng là lề trái của trang.
%Tóm tắt luận án phải phản ảnh trung thực kết cấu, bố cục và nội dung của luận án, phải ghi đầy đủ toàn văn kết luận của luận án.
%Mẫu trình bày trang bìa của tóm tắt luận văn (phụ lục 1).

% Các công trình liên quan
\chapter{Trình bày báo cáo}
\label{Chapter2}

\section{Quy định chung}

Báo cáo phải được trình bày ngắn gọn, rõ ràng, mạch lạc, sạch sẽ, không được tẩy xóa, có đánh số trang, đánh số bảng biểu, hình vẽ, đồ thị. 

Nội dung báo cáo được phân thành các chương. Số thứ tự của các chương, mục được đánh số bằng hệ thống số Ả-rập, không dùng số La mã. Các mục và tiểu mục được đánh số bằng các nhóm hai hoặc ba chữ số, cách nhau một dấu chấm: số thứ nhất chỉ số chương, chỉ số thứ hai chỉ số mục, số thứ ba chỉ số tiểu mục.


%Báo cáo cần dùng LaTEX để viết và trình bày theo mẫu đã được cung cấp.

 Báo cáo trình bày sử dụng khổ giấy với việc canh lề như sau: Lề trên 3 cm, lề dưới 2,5 cm, lề trái 3 cm, lề phải 2 cm. Đánh số trang ở giữa bên dưới. Đánh số trang ở giữa bên dưới.

Font chữ dùng trong báo cáo (Times New Roman) với kích cỡ (size) 13-14pt, sử dụng chế độ dãn dòng (line spacing) chế độ 1.5 lines.

%Các bảng biểu trình bày theo chiều ngang khổ giấy thì đầu bảng là lề trái của trang. 


\section{Bố cục của báo cáo}

Nội dung của báo cáo tối thiểu 50 trang khổ A4 và không nên vượt quá 100 trang (không kể các trang bìa, lời cám ơn, mục lục, tài liệu tham khảo \ldots) theo trình tự như sau:

\begin{itemize}
\item MỞ ĐẦU (thường đặt tên là ``Giới thiệu''): Trình bày lý do chọn đề tài, mục đích, đối tượng và phạm vi nghiên cứu.
Mô tả bài toán mà đề tài giải quyết.
Bài toán này có gì hay?
Tại sao lại cần giải quyết bài toán này?
Bài toán này có gì khó?
Có những hướng nào để giải quyết bài toán này?
Những hướng giải quyết trước đây có những vấn đề gì chưa giải quyết được?
Các câu hỏi nghiên cứu mà đề tài trả lời hoặc những vấn đề mà đề tài sẽ giải quyết.
Các đóng góp của đề tài.

\item TỔNG QUAN (thường đặt tên là ``Các công trình liên quan''): Phân tích đánh giá các hướng nghiên cứu đã có của các tác giả trong và ngoài nước liên quan đến đề tài; nêu những vấn đề còn tồn tại (những vấn đề nào mà các công trình khác chưa giải quyết được); chỉ ra những vấn đề mà đề tài cần tập trung, nghiên cứu giải quyết.

\item NGHIÊN CỨU THỰC NGHIỆM HOẶC LÝ THUYẾT (thường đặt tên là ``Phương pháp đề xuất''): Trình bày cơ sở lý thuyết, lý luận, giả thiết khoa học và phương pháp nghiên cứu đã được sử dụng trong đề tài.

Nếu đề xuất hướng giải quyết mới, mô hình mới thì cần mô tả chi tiết cách giải quyết của mình (chi tiết tới mức người khác có thể dựa vào phần này mà cài đặt lại được đúng hoàn toàn phương pháp của mình đề ra).

\item TRÌNH BÀY, ĐÁNH GIÁ BÀN LUẬN VỀ CÁC KẾT QUẢ (thường đặt tên là ``Kết quả thí nghiệm''): Mô tả các kết quả nghiên cứu khoa học hoặc kết quả thực nghiệm.
Đối với  đề tài ứng dụng có kết quả là sản phẩm phần mềm phải có hồ sơ thiết kế, cài đặt,\ldots theo một trong các mô hình đã học (UML,\ldots).

Thông thường cần mô tả môi trường thí nghiệm trước như sử dụng dữ liệu nào, dùng độ đo nào để đánh giá, môi trường chạy thí nghiệm (cấu hình máy nếu cần phân tích thông tin về thời gian chạy thực nghiệm). Sau đó, nêu kết quả thực nghiệm, bàn luận và giải thích kết quả.

\item KẾT LUẬN VÀ HƯỚNG PHÁT TRIỂN (thường đặt tên là ``Kết luận''): Trình bày những kết quả đạt được, những đóng góp mới và những đề xuất mới, kiến nghị về những hướng nghiên cứu tiếp theo.

\item DANH MỤC TÀI LIỆU THAM KHẢO: Chỉ bao gồm các tài liệu được trích dẫn, sử dụng và đề cập tới để bàn luận trong báo cáo.
Phần này các bạn chuẩn bị 1 file BIB để lưu các tài liệu trích dẫn.
Khi các bạn trích dẫn một tài liệu nào đó, LaTeX sẽ tự động thêm vào danh mục tài liệu tham khảo giúp các bạn.
Các bạn xem hướng dẫn cách trích dẫn ở chương sau.

\item PHỤ LỤC: Phần này bao gồm nội dung cần thiết nhằm minh họa hoặc hỗ trợ cho nội dung báo cáo như số liệu, mẫu biểu, tranh ảnh,\ldots Phụ lục không được dày hơn phần chính của báo cáo.
Nếu có công trình công bố thì để vào phần phụ lục này.
\end{itemize}

\section{Bảng biểu, hình vẽ, phương trình}

%Những qui định dưới này các bạn có thể bỏ qua hoặc đọc để hiểu thêm.
%Những định dạng này LaTeX đều tự động giúp các bạn.
%Các bạn xem hướng dẫn chi tiết hơn ở chương sau.

Việc đánh số bảng biểu, hình vẽ, phương trình phải gắn với số chương; ví dụ hình 3.4 có nghĩa là hình thứ 4 trong Chương 3.
Mọi đồ thị, bảng biểu, hình vẽ lấy từ các nguồn khác phải được trích dẫn đầy đủ.

\subsection{Bảng biểu, hình vẽ}



Đầu đề của bảng biểu ghi phía trên bảng, đầu đề của hình vẽ ghi phía dưới hình.

Thông thường, những bảng ngắn và đồ thị phải đi liền với phần nội dung đề cập tới các bảng và đồ thị này ở lần thứ nhất.
Các bảng dài có thể để ở những trang riêng nhưng cũng phải tiếp theo ngay phần nội dung đề cập tới bảng này ở lần đầu tiên.
Các bảng rộng vẫn nên trình bày theo chiều đứng dài 297mm của trang giấy, chiều rộng của trang giấy có thể hơn 210mm.
Chú ý gấp trang giấy sao cho số và đầu đề của hình vẽ hoặc bảng vẫn có thể nhìn thấy ngay mà không cần mở rộng tờ giấy.
Tuy nhiên hạn chế sử dụng các bảng quá rộng này.

Đối với những trang giấy có chiều đứng hơn 297mm (bản đồ, bản vẽ,\ldots) thì có thể để trong một phong bì cứng đính bên trong bìa sau của báo cáo.

Các hình vẽ phải sạch sẽ bằng mực đen để có thể sao chụp lại; có đánh số và ghi đầy đủ đầu đề, cỡ chữ phải bằng cỡ chữ sử dụng trong báo cáo.

Khi đề cập đến các bảng biểu và hình vẽ phải nêu rõ số của hình và bảng biểu đó, ví dụ ``... được nêu trong Bảng 4.1'' hoặc ``xem Hình 3.2'' mà không được viết ``… được nêu trong bảng dưới đây'' hoặc ``trong đồ thị của X và Y sau''.

\subsection{Phương trình toán học}

Việc trình bày phương trình toán học trên một dòng đơn hoặc dòng kép tùy ý, tuy nhiên phải thống nhất trong toàn báo cáo.

Khi ký hiệu xuất hiện lần đầu tiên thì phải giải thích và đơn vị tính phải đi kèm ngay trong phương trình có ký hiệu đó.
Nếu cần thiết, danh mục của tất cả các ký hiệu, chữ viết tắt và nghĩa của chúng cần được liệt kê và để ở phần đầu của báo cáo.

Tất cả các phương trình cần được đánh số và để trong ngoặc đơn đặt bên phía lề phải.
Nếu một nhóm phương trình mang cùng một số thì những số này cũng được để trong ngoặc, hoặc mỗi phương trình trong nhóm phương trình (5.1) có thể được đánh số là (5.1.1), (5.1.2), (5.1.3).

\section{Viết tắt}

\textbf{Không lạm dụng việc viết tắt} trong báo cáo.
Chỉ viết tắt những từ, cụm từ hoặc thuật ngữ được sử dụng nhiều lần trong báo cáo.
Không viết tắt những cụm từ  dài, những mệnh đề; không viết tắt những cụm từ ít xuất hiện trong báo cáo.
Nếu cần viết tắt những từ thuật ngữ, tên các cơ quan, tổ chức,\ldots thì được viết tắt sau lần viết thứ nhất có kèm theo chữ viết tắt trong ngoặc đơn.
Nếu báo cáo có nhiều chữ viết tắt thì phải có bảng danh mục các chữ viết tắt (xếp theo thứ tự ABC) ở phần đầu báo cáo.

Nhắc lại: \textbf{không lạm dụng việc viết tắt} trong báo cáo.
Khi các bạn sử dụng từ viết tắt, người đọc sẽ phải lật lại những phần đã đọc, để tìm lại xem từ viết tắt đó nghĩa là gì.
Việc này sẽ làm chậm tốc độ đọc và sẽ khiến người đọc khó theo dõi báo cáo của bạn hơn.
Nếu có thể, hạn chế hoàn toàn việc dùng viết tắt.

\section{Tài liệu tham khảo và cách trích dẫn}

Mọi ý kiến, khái niệm có ý nghĩa, mang tính chất gợi ý không phải của riêng tác giả và mọi tham khảo khác phải được trích dẫn và chỉ ra nguồn trong danh mục Tài liệu tham khảo của báo cáo. Nguồn được trích dẫn phải được liệt kê chính xác trong danh mục Tài liệu tham khảo.

Việc trích dẫn, tham khảo chủ yếu nhằm thừa nhận nguồn của những ý tưởng có giá trị giúp người đọc theo được mạch suy nghĩ của tác giả, không làm trở ngại việc đọc.


Không trích dẫn những kiến thức phổ biến, mọi người đều biết cũng như không làm báo cáo nặng nề với những tham khảo trích dẫn.

Nếu không có điều kiện tiếp cận được một tài liệu gốc mà phải trích dẫn thông qua một tài liệu khác thì phải nêu ra trích dẫn này, đồng thời tài liệu gốc đó không được liệt kê trong danh mục tài liệu tham khảo của báo cáo.
% \chapter{Phương pháp dựa trên luật}
\label{Chapter2}

\section{Giới Thiệu}

Hiện nay các bài toán liên quan đến dự đoán liên kết đồ thị tri thức lớn rất được quan tâm có khoảng bốn nhánh nghiên cứu chính như được nhắc đến trong nghiên cứu \cite{ampligraph} một trong số đó là phương pháp dựa trên luật logic. Với mong muốn được tiếp cận các phương pháp từ đơn giản đến phức tạp nên chúng tôi chọn phương pháp này làm chủ để chính cho các báo cáo và nghiên cứu trong chương này. Với phương pháp này đưa ra một xếp hạng \(k\) ứng viên với một số điểm nhất định biểu thị cho sự chắc chán của dự đoán nó phù hợp với các hệ thống gợi ý(recommender system).

\section{Các công trình liên quan}
Hầu hết các nghiên cứu hiện tại về việc dự đoán liên kết của đồ thị tri thức đều liên quan đến các phương pháp tiếp cận tập trung vào khái niệm nhúng một đồ thị đã cho trong một không gian vectơ có số chiều thấp. Ngược lại với các tiếp cận này là một phương pháp đựa trên luật được nghiên cứu trong \cite{burl}.Thuật toán cốt lõi của nó dựa trên lấy mẫu một luật bất kỳ, sau đó khái quát  thành các quy tắc Horn\cite{wiki:Horn} Tiếp đó dùng thống kê để tính độ tin cậy của các luật được khái quát. Khi dự đoán một liên kết mới (cạnh mới) của đồ thị chúng ta dự đoán một đỉnh có cạnh nối với một quan hệ cụ thể (label) với đỉnh còn lại hay không.Cũng đã có rất nhiều phương pháp được nghiên cứu, đề xuất để học các các luật trong đồ thị chẳng hạn như trong   RuDiK \cite{dettmers2018convolutional}, AMIE\cite{galarraga2015fast}, RuleN\cite{meilicke2018fine}. Trong phương pháp RuDiK học các luật có thể học cả quy tắc positive và negative. Trong khi phương pháp của chúng tôi chỉ học các quy tắc positive. RuDiK tìm ra một bộ quy tắc nhỏ bao gồm phần lớn các ví dụ là đúng và ít sai sót nhất có thể. Điều này khác với mục tiêu của chúng tôi, chúng tôi cố gắng tìm hiểu mọi quy tắc khả thi có thể sau đó tạo xếp hạng \(k\) ứng viên tiềm năng với một độ tin cậy nhất định đuọc đo trên tập huấn luyện.

Phương pháp của chúng tôi phần lớn dựa vào phương pháp Anytime Bottom-Up Rule Learning for Knowledge Graph Completion \cite{meilicke2019anytime} mà sau đây chúng tôi gọi là \textbf{AnyBURL}. Vấn đề tồn đọng lại ở mô hình này khi có một cạnh mới hay một tri thức mới được thêm vào đồ thị sẽ phải đào tạo lại toàn bộ mô hình. Chúng tôi giải quyết vẫn đề này theo hai chiến lược offline-to-online tức là khi thêm vào đồ thị tập hợp các cạnh thì mới thực hiện lại quá trình đào tạo lại một phần của đồ thị và chiến lược thứ 2 là online-to-online  khi thêm một cạnh mới sẽ thực hiện đào tạo lại ngay một phần có liên quan tới cạnh vừa thêm vào.
\section{Phương pháp đề xuất}
Trong phần này chúng tôi mô tả lại cách mô hình hóa lại bài toán theo phương pháp dựa trên luật AnyBURL,thuật toán lấy mẫu luật (đường đẫn) và thuật toán khái quát hóa một luật để lưu trữ trở thành tri thức của mô hình. Cùng với những cải tiến của chúng tôi trong quá trình đào tạo khi có một tri thức mới được thêm vào đồ thị(thêm cạnh).
\subsection{Horn rule}
Trong logic toán học, một công thức nguyên tử - \textbf{atomic formula}\cite{wiki:Atomic} (còn được gọi đơn giản là một nguyên tử-\textbf{atom}) là một công thức không có cấu trúc mệnh đề, nghĩa là một công thức không chứa các liên kết logic (\(\vee, ~ \wedge\)) hoặc tương đương (\(\Leftrightarrow\)) là một công thức không có các mẫu con nghiêm ngặt (tức là atom không thể chia nhỏ ra thành các atom con nữa). Do đó, các nguyên tử là công thức đơn giản nhất để hình thành luật của logic. Các công thức hợp được hình thành bằng cách kết hợp các công thức nguyên tử bằng cách sử dụng các liên kết logic.

Một \textbf{literal}\cite{wiki:Literal} là một công thức nguyên tử (nguyên tử) hoặc phủ định của nó. Định nghĩa chủ yếu xuất hiện trong lý thuyết logic cổ điển. \textbf{Literal} có thể được chia thành hai loại: Một \textbf{positive literal} chỉ là một nguyên tử (ví dụ: \(x\)). Một \textbf{negative literal} là phủ định của một nguyên tử (ví dụ: \(\neg x\)). Sự phân chia của \textbf{literal} là \textbf{positive literal} hay \textbf{negative literal} tùy thuộc vào việc \textbf{literal} được định nghĩa.

Một mệnh đề (clause) là một literal hoặc nối rời của hai hoặc nhiều literal. Ở dạng \textbf{Horn} một mệnh đề có nhiều nhất một positive literal. Lưu ý: Không phải mọi công thức trong logic mệnh đề đều có thể đưa về dạng Horn.Mệnh đề xác định không có literal đôi khi được gọi là mệnh đề đơn vị (unit clause) và một mệnh đề đơn vị không có biến đôi khi được gọi là \textit{facts}\cite{wiki:Horn}.Một công thức nguyên tử được gọi là \textit{ground} hoặc \textit{ground atoms} nếu nó được xây dựng hoàn toàn từ các mệnh đề đơn vị; tất cả các \textit{ground atoms} có thể ghép lại từ một tập hợp hàm và các ký hiệu vị từ nhất định tạo nên cơ sở Herbrand cho các bộ ký hiệu này\cite{wiki:Term}.

\subsection{Định nghĩa đồ thị  tri thức} \label{kg}
Một đồ thị tri thức \(\mathbb{G}\) được định nghĩa trên một bộ từ vựng \(\langle \mathbb{C}, \mathbb{R} \rangle\) trong đó \(\mathbb{C}\) là tập hợp các hằng số và \(\mathbb{R}\) là tập hợp các vị từ nhị phân.Khi đó, \(\mathbb{G} = \{r (a, b) \mid r \in \mathbb{R}, a, b \in \mathbb{C}\}\) là tập hợp các \textit{ground atoms} hoặc \textit{facts}. Một vị từ nhị phân được gọi là quan hệ và hằng số (hoặc hằng số được đề cập đến) được gọi là thực thể (entity) tương ứng với một dòng dữ liệu trong tập huấn luyện. Sau đây chúng tôi sử dụng các chữ cái viết thường cho các hằng và chữ in hoa cho các biến cho các thảo luận dưới đây. Vì chúng ta không học các quy tắc Horn tùy ý, và chỉ học đối với loại quy tắc nào có thể được khái quát hóa như được thảo luận dưới đây.

Chúng ta định nghĩa một quy tắc là \(h(c_0, c_n) \gets b_1(c_0, c_1) ,\dots ,b_n(c_{n}, c_{n + 1})\) là một đường dẫn ground atoms có chiều dài \(n\). Trong đó \(h(\dots)\) được gọi là \textit{head atoms} và \( b_1(c_0, c_1) ,\dots ,b_n(c_{n}, c_{n + 1})\) được gọi là \textit{body atoms}. Chúng tôi sẽ phân biệt dưới đây ba loại quy tắc mà chúng tôi gọi là: \textit{quy tắc nhị phân} \((\mathbf{B})\) là quy tắc trong head atoms chứa 2 biến, quy tắc đơn nguyên kết thúc bằng một đỉnh treo  và atom này chỉ chứa biến không chứa hằng số\((\mathbf{U_d})\) và head atoms chỉ chứa 1 biến. Còn quy tắc đơn nguyên kết thúc bằng một atom \((\mathbf{U_c)}\) và head atoms cũng chỉ chứa 1 biến.\((\mathbf{U_c)}\)  có thể là một đỉnh treo tới một hằng số bất kì nếu hằng số này trùng mới hằng số trong head atom thì tạo thành một đường đẫn có chu trình.

\[B \hspace{3.7cm} h(A_0,A_n) \gets  \bigwedge^n_{i=1} b_i(A_{i-1}, A_i)\]
\[U_d \hspace{3.8cm} h(A_0,c) \gets  \bigwedge^n_{i=1} b_i(A_{i-1}, A_i)\]
\[U_c \hspace{1cm} h(A_0,c) \gets  \bigwedge^{n-1}_{i=1} b_i(A_{i-1}, A_i) \wedge b_n(A_{n-1}, c^{\prime})\]

Chúng tôi gọi các quy tắc của các loại này là quy tắc đường đi (path rules), bởi vì các body atoms (phần sau đấu \(\gets\)) tạo thành một đường đi. Lưu ý rằng nó cũng bao gồm các biến thể quy tắc với các biến được đảo ngược trong các nguyên tử: được đưa ra trong đồ thị tri thức \(\mathbb{G}\), đường dẫn có độ dài \(n\) là một chuỗi gồm \(n\) bộ ba \(p_i (c_i, c_i + 1)\) với \(p_i (c_i, c_i + 1) \in \mathbb{G}\) hoặc \(p_i (c_i + 1, c_i) \in \mathbb{G}\) với \(0 \geq i \leq n\). Các mẫu quy tắc trừu tượng (abstract rule patterns) được cho ở trên có độ dài \(n\) vì body atoms của chúng có thể được khởi tạo thành một đường dẫn có độ dài \(n\).

Ngoài ra Quy tắc \(B\) và quy tắc \(U_c\) cũng được gọi là quy tắc kết nối kín. Chúng có thể được học bởi hệ thống khai thác AMIE được mô tả trong \cite{AMIE,galarraga2015fast}. Quy tắc \(U_d\) là quy tắc không đóng hay đường đi không tạo thành chu trình vì \(A_n\) là biến chỉ xuất hiện một lần. Ví dụ:
\[
\begin{matrix}
\textit{speaks}(X, Y ) & \gets & \textit{lives}(X, Y) & \quad (1) \\
\textit{lives\_in\_city}(X, Y ) & \gets & \textit{lives}(X, A),\textit{within}(Y, A)  & \quad  (2) \\
\textit{gen}(X, female) & \gets & \textit{married}(X, A), \textit{gen}(A, male)  & \quad  (3) \\
\textit{profession}(X, actor) &  \gets & \textit{acted\_in}(X, A)  & \quad (4)
\end{matrix}
\]
Quy tắc (1) là quy tắc \textbf{B}(quy tắc nhị phân) quy tắc này nói rằng nếu một người (thực thể) \(X\) nói nguôn ngữ \(Y\) nếu người \(X\) sống  ở đất nước \(Y\).Rõ ràng quy tắc này là một quy tắc khái quát miễn khi nào thực thể \(X\) có cạnh nối với thực thể \(Y\) với nhãn là \textit{lives} thì có thể kết thêm 1 cạnh với nhãn \textit{speaks} giữa \(X\) và \(Y\). Quy tắc (2), (3) điều là quy tắc \(U_c\) ,quy tắc (2) nói rằng người \(X\) sống ở thành phố \(Y\) nếu người \(X\) sống ở quốc gia \(A\) và thành phố \(Y\) nằm trong quốc gia \(A\), quy tắc (3) nói rằng nếu một người \(X\) là nữ nếu họ kết hôn với một người \(A\) và người \(A\) có giới tính nam. Ở quy tắc (3) không tạo thành chu trình trên đồ thị như quy tắc (2) đỉnh (Y) lặp lại  ở \textit{head atom} và đỉnh cuối cùng trong \textit{body atoms}. Quy tắc (4) là quy tắc \(U_d\) nói rằng người \(X\) là một điễn viên nếu người \(X\) đóng trong một bộ phim \(A\).

Tất cả các quy tắc được xem xét sẽ được lọc lại đựa trên điểm được gọi là độ tin cậy của quy tắc là được đo trên tập dữ liệu huấn luyện. Độ tin cậy này được đo bằng tỷ lệ body atoms dẫn đến head atoms chia cho tất cả các đường đãn chứa body atoms.Ví dụ khi ta có quy tắc sau:
\(\textit{gen}(X, female) \gets \textit{married}(X, A), \textit{gen}(A, male) \). Khi đó chúng ta thực hiện đếm tất cả các cặp thực thể có quan hệ  \(\textit{married}(X, A), \textit{gen}(A, male) \) được gọi là số đường dẫn chứa body atoms, sau đó thực hiện đếm tất cả các  thực thể thỏa quan hệ \(\textit{gen}(X, female) \gets \textit{married}(X, A), \textit{gen}(A, male) \) được gọi là số body atoms dẫn đến head atoms. Sau đó chia số body atoms dẫn đến head atoms cho  đường dẫn chứa body atoms được gọi là độ tin cậy của quy tắc.
\subsection{Thuật toán} \label{algorithm2}
Trong phần này chúng tôi mô tả lại thuật toán chính của phương pháp AnyBURL nó cũng được mô tả trong \cite{burl} cũng như hai thuật toán mở rộng của chúng tôi để giải quyết vấn đề khi đồ thị được thêm một hoặc một lượng tri thức mới (thêm cạnh). Ngoài ra chúng tôi cũng mô tả sơ lược lại cách khởi tạo một luật cũng như cách thức tính toán độ tin cậy bằng cách lấy mẫu trên tập huấn luyện và vấn đề độ tin cậy khi dự đoán một luật khi tính toán độ tin cậy bằng việc lấy mẫu.
\subsubsection{Thuật toán 1 AnyBURL}
\begin{algorithm}
\caption{Anytime Bottom-up Rule Learning}\label{euclid}
\begin{algorithmic}[1]
\Procedure{AnyBURL($\mathbb{G}$, s, sat, Q, ts)}{}
\State $\textit{n} = \text{2}$
\State $R = \emptyset$
\Loop
\State $R_s = \emptyset$
\State $start = currentTime()$
\Repeat
\State $p = samplePath(\mathbb{G}, n)$
\State $R_p = generateRules(p)$
\For {$r \in R_p$}
\State $score(r, s)$
\If {$Q(r)$}
	\State $R_s = R_s \cup \{r\}$
\EndIf
\EndFor
\Until {$currentTime() > start + ts$}
\State $R^{\prime}_s = R_s \cap R$
\If {$ \mid R^{\prime} \mid / \mid R \mid > SAT$}
	\State $n = n + 1$
\EndIf
\State $R = R_s \cap R$
\EndLoop
\Return R
\EndProcedure
\end{algorithmic}
\end{algorithm}

Đầu vào của thuật toán \(\mathbb{G}, S, SAT, Q, TS\). Đầu ra là tập hợp \(R\) các luật học được. Trong đó \(\mathbb{G}\) là một đồ thị tri thức được cho từ tập dữ liệu đào tạo. \(S\) là tham số cho biết kích thước của một lần lấy mẫu trên dữ liệu đào tạo để tính toán độ tin cậy. \(SAT\) cho biết độ bão hòa(saturation) của các luật được sinh ra trong 1 lần lặp độ bão hòa này được tính bằng số luật \textbf{mới} học được ở lần lặp hiện tại so số sô luật đã học được. Nếu nhỏ hơn độ bão hòa thì chúng tôi cho rằng vẫn còn tiềm năng để khai thác các luật với độ dài \(n\). ngược lại chúng tôi tăng độ dài của luật sau đó tiếp tục khai thác. \(Q\) là một ngưỡng để xác định xem luật mới được sinh ra có được thêm vào kết quả trả về hay không. Còn \(TS\) cho biết thời gian học của thuật toán. Chúng tôi bắt đầu với \(n\) bằng \(2\) tức là các luật có độ dài đường đẫn bằng 2 vì trong path rule yêu cầu ít nhất 1 literal trong head atom và 1 trong body atoms. Ở phần lấy mẫu 1 luật(\textit{samplePath}) chỉ đơn giản là chúng ta chọn 1 đỉnh bất kì trong đồ thị duyệt qua tất cả các đường đẫn từ đỉnh đó đi qua \(n\) đỉnh khác, sau đó chọn ngẫu nhiên 1 đường đẫn trong số các trường đẫn duyệt được.

\subsubsection{Thuật toán 2 tạo 1 luật}
\begin{algorithm}
\caption{Generate Rules(p)}\label{euclid}
\begin{algorithmic}[1]
\Procedure{generate\_rules(p)}{}
\State $\textit{generalizations} = \emptyset$
\State $is\_binary\_rule = random.choices([true,false])$
\If {$is\_binary\_rule$}
	\State $replace\_all\_head\_by\_variables(p)$
	\State $replace\_all\_tail\_by\_variables(p)$
	\State $add(generalizations, p)$
\Else:
    \State $replace\_all\_head\_by\_variables(p)$
    \State $add(generalizations, p)$
    \State $replace\_all\_tail\_by\_variables(p)$
	\State $add(generalizations, p)$
\EndIf
\Return $generalizations$
\EndProcedure
\end{algorithmic}
\end{algorithm}

Ở thuật toán này chúng tôi thay các hằng số vào các head và tail trong toàn bộ path rule  của luật được lấy mẫu ở bước trước nếu luật cần học là luật nhị phân ngược lại chúng tôi chỉ thay hoặc head hoặc tail rồi thêm vào luật trả về sau đó chúng tôi lấy mẫu trên tập huấn luyện 1 tập hợp các luật sau đó tính toán độ tin cây như được mô tả trong phần \hyperref[kg]{2.3.2}. Để giảm chi phí tính toán chúng tôi chọn cách lấy mẫu trên tập huấn luyện để tính toán. Khi đưa ra dự đoán các ứng cử viên của một luật chúng tôi sẽ tính toán lại bằng cách thêm vào một lượng biểu diẽn số luật bị sai mà chúng tôi chưa nhìn thấy trong quá trình lấy mẫu để tính toán độ tin cậy. Dối với mô hình của chúng thôi sau khi thử nghiệm tham số trong khoảng \([5, 10]\) cho kết quả tốt nhất.

\subsubsection{Thuật toán 3 học offline-to-online}
\begin{algorithm}
\caption{AnyBURL Learning batch size}\label{euclid}
\begin{algorithmic}[1]
\Procedure{AnyBURLbatch($\mathbb{G}$, s, sat, Q, ts, batch\_edge)}{}
\State $is\_connected = add(\mathbb{G}, batch\_edge)$
\If {$is\_connected$}
	\State  $ G^{\prime} = \mathbb{G} \oplus batch\_edge$
\Else
    \State  $ G^{\prime} = batch\_edge$
\EndIf
\State $\textit{n} = \text{2}$
\State $R = \emptyset$
\Loop
\State $R_s = \emptyset$
\State $start = currentTime()$
\Repeat
\State $p = samplePath(G^{\prime}, n)$
\State $R_p = generateRules(p)$
\For {$r \in R_p$}
\State $score(r, s)$
\If {$Q(r)$}
	\State $R_s = R_s \cup \{r\}$
\EndIf
\EndFor
\Until {$currentTime() > start + ts$}
\State $R^{\prime}_s = R_s \cap R$
\If {$ \mid R^{\prime} \mid / \mid R \mid > SAT$}
	\State $n = n + 1$
\EndIf
\State $R = R_s \cap R$
\EndLoop
\Return R
\EndProcedure
\end{algorithmic}
\end{algorithm}

Thuật toán này là phần bổ xung của chúng tôi để tránh việc phải đào tạo lại toàn bộ mô hình khi có một lượng tri thức mới được thêm vào đồ thị. Khi thêm vào đồ thị chúng tôi kiểm trả xem phần tri thức mới có kết nối với tri thức cũ hay không (tính liên thông) nếu có chúng tôi thực hiện phép toán \(\oplus\) lấy  tất cả các phần trong \(batch\_edge\) thêm với 1 phần liên thông với những cạnh liên thông với đồ thị với dộ dài là \(5\), Nếu không chúng tôi lấy tất cả các phần trong \(batch\_edge\) sau đó thực hiện lại các bước như thuật toán Anytime Bottom-up Rule Learning.

\subsubsection{Thuật toán 4 học online-to-online}
\begin{algorithm}
\caption{AnyBURL Learning batch size}\label{euclid}
\begin{algorithmic}[1]
\Procedure{AnyBURLbatch($\mathbb{G}$, s, sat, Q, ts, edge)}{}
\State $is\_connected = add(\mathbb{G}, edge)$
\State $R = \emptyset$
\If {$is\_connected$}
    \State $\textit{n} = \text{2}$
    \State $R_s = \emptyset$
    \Repeat
        \State $p = samplePath(edge, n)$
        \State $R_p = generateRules(p)$
        \For {$r \in R_p$}
            \State $score(r, s)$
            \If {$Q(r)$}
        	    \State $R_s = R_s \cup \{r\}$
            \EndIf
        \EndFor
    \Until {$currentTime() > start + ts$}
    \State $R^{\prime}_s = R_s \cap R$
    \If {$ \mid R^{\prime} \mid / \mid R \mid > SAT$}
    	\State $n = n + 1$
    \EndIf
    \State $R = R_s \cap R$
\State \Return R
\EndIf
\EndProcedure
\end{algorithmic}
\end{algorithm}

Thuật toán này là một phần bổ xung cho thuật toán 3 ở trên. Sở dĩ chúng tôi gọi là online-to-online là vì khi có một cạnh mới(tri thức mới) được thêm vào đồ thị chúng tôi sẽ thực hiện việc học ngay tức khắc trên các path rule liên quan tới cạnh đó không giống như ở thuật toán 3 khi có đủ 1 lượng tri thức mới được thêm vào.

% Nội dung báo cáo được phân thành các chương. Số thứ tự của các chương, mục được đánh số bằng hệ thống số Ả-rập, không dùng số La mã. Các mục và tiểu mục được đánh số bằng các nhóm hai hoặc ba chữ số, cách nhau một dấu chấm: số thứ nhất chỉ số chương, chỉ số thứ hai chỉ số mục, số thứ ba chỉ số tiểu mục.

% %Báo cáo cần dùng LaTEX để viết và trình bày theo mẫu đã được cung cấp.

%  Báo cáo trình bày sử dụng khổ giấy với việc canh lề như sau: Lề trên 3 cm, lề dưới 2,5 cm, lề trái 3 cm, lề phải 2 cm. Đánh số trang ở giữa bên dưới. Đánh số trang ở giữa bên dưới.

% Font chữ dùng trong báo cáo (Times New Roman) với kích cỡ (size) 13-14pt, sử dụng chế độ dãn dòng (line spacing) chế độ 1.5 lines.

%Các bảng biểu trình bày theo chiều ngang khổ giấy thì đầu bảng là lề trái của trang.

\section{Kết quả thí nghiệm}
Trong phần này chúng tôi mô tả lại các bộ dữ liệu mà chúng tôi đùng để thực nghiệm đánh giá phương pháp của chúng tôi cùng với so sánh với các kết quả khác của các công trình nổi bật khác được báo cáo trong \hyperref[tab:tab1]{bảng 2.1}. Ngoài ra chúng tôi cũng cố gắng đánh giá hai đề xuất của chúng tôi trong việc thêm một lượng tri thức mới vào đồ thị bằng cách chúng tôi xem tập test là một lượng tri thức mới cần thêm vào và dùng tập vadidation để đánh giá lại phương pháp của chúng tôi. Kết quả chi tiết được báo cáo ở \hyperref[tab:tab1]{bảng 2.2}
\subsection{Các tập dữ liệu huấn luyện} \label{datasets}
Trong thí nghiệm của chúng tôi, chúng tôi thực hiện trên bốn tập dữ liệu phổ biến là FB15k, FB15-237, WN18 và WN18RR.Các bộ dữ liệu này là tập hợp các bộ ba (triple) \(\langle head, relation, tail \rangle\) biểu thị cho thực thể đầu có một mối quan hệ với thực thể cuối.

Bộ dữ liệu FB15k: Bộ dữ liệu này được tạo bởi nhóm nghiên cứu A. Bordes, N. Usunier, A. Garcia-Duran, J. Weston, and O. Yakhnenko \cite{bordes2013translating} bằng cách trích xuất từ bộ dữ liệu Wikilinks database \footnote{https://code.google.com/archive/p/wiki-links/}.Wikilinks database thu thập các siêu liên kết(hyperlinks) đến Wikipedia gồm 40 triệu lượt đề cập trên 3 triệu thực thể, họ trích xuất tất cả các dữ kiện liên quan đến một thực thể nhất định có hơn 100 lần được đề cập đến bởi các tài liệu khác cùng với tất cả các dữ kiện liên quan đến thực thể đó (bao gồm cả những thực thể con được nhắc đến trong tài liệu Wikipedia đó), ngoại trừ những thonng tin như: ngày tháng, danh từ riêng, v.v ... Họ cũng chuyển đổi các đỉnh có bậc \(n\) được biểu diễn thành các nhóm các cạnh nhị phân tức là liệt kê các cạnh và quan hệ của mọi đỉnh. Tập dữ liệu được chia ngẫu nhiên thành 3 tập: tập training với 1345 relations, 14834 head entities và 14903 tail entities, tập test gồm 916 relations, 11886 head entities, và 11285 tail entities, tập vadiation gồm 961 relations, 12297 head entities, và 11825 tail entities.

Bộ dữ liệu FB15k-237: là một tập hợp con của FB15k được xây dựng bởi Toutanova và Chen \cite{toutanova2015observed} lấy cảm hứng từ quan sát rằng FB15k bao gồm dữ liệu thử nghiệm được các mô hình nhìn thấy tại thời điểm đào tạo(test leekage). Trong FB15k, vấn đề này là do sự hiện diện của các quan hệ gần giống nhau hoặc nghịch đảo của nhau.FB15k-237 được xây dựng để trở thành một tập dữ liệu thách thức hơn: các tác giả đã chọn các dữ kiện liên quan đến 401 quan hệ xuất hiện nhiều nhất và loại bỏ tất cả các quan hệ tương đương hoặc nghịch đảo. Họ cũng đảm bảo rằng không có thực thể nào được kết nối trong tập huấn luyện cũng được liên kết trực tiếp trong tập test và validation.Tập training gồm 237 relations, 13781 head entities, và 13379 tail entities, tập test gồm 223 relations, 7652 head entities, và 5804 tail entities, tập vadition gồm 224 relations, 8171 head entities, and 6376 tail entities.

Bộ dữ liệu WN18: được giới thiệu bởi các tác giả của TransE \cite{bordes2013translating}, được trích xuất từ WordNet\footnote{https://wordnet.princeton.edu/}, một bản thể học ngôn ngữ KG có nghĩa là cung cấp một từ điển/từ đồng nghĩa để hỗ trợ NLP và phân tích văn bản tự động. Trong WordNet, các thực thể tương ứng với các tập hợp (\textit{word senses}) và các quan hệ đại diện cho các kết nối từ vựng của chúng (ví dụ: “hypernym”). Để xây dựng WN18, các tác giả đã sử dụng WordNet làm điểm bắt đầu và sau đó lặp đi lặp lại lọc ra các thực thể và mối quan hệ với quá ít lần được đề cập. Tập dữ liệu được chia ngẫu nhiên thành 3 tập: tập training với 18 relations, 40504 head entities, và 40551 tail entities, tập test gồm 18 relations, 4262 head entities, and 4338 tail entities, tập vadiation gồm 18 relations, 4349 head entities, and 4263 tail entities.

Bộ dữ liệu WN18RR: là một tập hợp con của WN18 được xây dựng bởi DeŠmers et al.\cite{dettmers2017convolutional}, cũng là người giải quyết vấn đề rò rỉ thử nghiệm (test leakage)trong WN18. Để giải quyết vấn đề đó, họ xây dựng tập dữ liệu WN18RR thách thức hơn nhiều bằng cách áp dụng một phương pháp tương tự đưược sử dụng cho FB15k-237 \cite{toutanova2015observed}. Training gồm 11 relations, 39610 head entities, và 31881 tail entities, tập test gồm 11 relations, 2958 head entities, và 2619 tail entities, tập vadition gồm 11 relations, 2851 head entities, and 2575 tail entities.

\subsection{Kết quả thực nghiệm}
Trong phần này chúng tôi mô tả lại các phương pháp đánh giá(độ do), môi trường thực hiện cũng như các tập dữ liệu mà chúng tôi sử dụng để dánh giá phương pháp của mình. Các phương pháp đánh giá(độ do) này cũng phổ biến nó được đánh giá cho hầu hết các mô hình dự đoán liên kết trên đồ thị. Chúng tôi tiến hành so sánh với bốn phương pháp nổi bật khác được báo cáo trong \cite{rossi2020knowledge}.
\subsubsection{Các độ đo}
\textit{Mean Rank} (MR). Đây là giá trị trung bình của rank thu được cho một dự đoán chính xác. Càng nhỏ thì mô hình càng chính xác:
\[MR = \frac{1}{\mid Q \mid} \sum_{q ~\in~ Q} rank(q) \]
Trong đó \(\mid Q \mid\) là độ lớn của tập hợp các câu hỏi bằng độ lớn của tập test hoặc vadidation. Khi dự đoán chúng tôi dự đoán cả head và tail cho một dòng tương ứng trong tập dữ liệu thử nghiệm. Ví dụ chún tôi sẽ dự đoán \(\langle ?,~ relation,~ tail \rangle\) và \(\langle head,~ relation,~ ?\rangle\) cho 1 dòng tương ứng. \(q\) thể hiện cho câu hỏi chúng tôi dự đoán và \(rank(q)\) thể hiện cho kết quả đúng của câu hỏi đứng ở vị trí thứ mấy trong xếp hạng của chúng tôi sau đó lấy trung bình rank của các dự đoán head và tail. Rõ ràng độ đo này nằm giữa \([1, \mid \text{số lượng các entity} \mid]\) do có tối da \(n\) cạnh nối 1 đỉnh tới \(n-1\) đỉnh còn lại và thêm cạnh nối tới chính đỉnh nó(cạnh khuyên). Và độ đo này đễ bị ảnh hưởng bởi nhiễu vì có những quan hệ có những thực thể được xếp hạng gần cuối. Để giải quyết vấn đề này nhóm chúng tôi và các nhóm nghiên cứu khác sử dụng thêm độ đo Mean Reciprocal Rank (MMR).

\textit{Mean Reciprocal Rank} (MMR). Đây là xếp hạng đối ứng trung bình, là nghịch đảo của giá trị trung bình của rank thu được cho một dự đoán chính xác ở trên. Và càng lớn thì mô hình càng chính xác. Do độ đo này lấy nghịch đảo của các rank nên tránh dược vấn đề nhiễu của độ đo MR ở trên.
\[MRR =\frac{1}{\mid Q \mid} \sum_{q~ \in ~Q} \frac{1}{rank(q)}\]

\textit{Hit@K} (H@K). Đó là tỷ lệ các dự đoán đúng mà rank nhỏ hơn hoặc bằng ngưỡng \(K\):
\[H@K = \frac{\mid {q ~\in ~Q~: rank(q) \leq K} \mid}{\mid Q \mid}\]
\subsubsection{Kết quả}
Như đã nói trước đây với mô hình dựa trên luật của chúng tôi hoàn toàn có thể thực hiện trên một laptop với cấu hình thông thường. Trong thí nghiệm của chúng tôi cấu hình máy để thực thi như sau: T480, core i5 8th Gen, ram 16Gb, 4 core 8 thread. Mã nguồn thực thi được viết bằng ngôn ngữ Python phiên bản 3.6 và dùng các hàm hỗ trợ có sẵn trong Python với không một thư viện bên thứ ba nào. Thí nghiệm được thực hiện với bốn tập dữ liệu phổ biến là FB15k, FB15-237, WN18 và WN18RR. Thông tin chi tiết các bộ dữ liệu này được mô tả ở phần \hyperref[datasets]{các tập dữ liệu huấn luyện}.

Như mô tả ở phần \hyperref[algorithm2]{thuật toán AnyBURL} thuật toán này sẽ học các luật được sinh ra trong một khoảng thời gian nhất định do người dùng cấu hình. Ở đây chúng tôi chọn cấu hình thời gian là 1000 giây tương đương khoảng 17 phút đào tạo, với độ bão hòa(SAT) \(0.85\), độ tin cậy Q \(0.05\), kích thước mẫu S (\(\frac{1}{10}~ \text{tập huấn luyện}\)). Với cấu hình như vậy mô hình phiên bản Python của chúng tối cho kết quả tương đương với phiên bản Java nhóm tác giả Meilicke, Christian et al. \cite{burl} với cấu hình tương tự nhưng thời gian training là 100 giây. Sự khác biệt về thời gian học tập ở đây chủ yếu là do hiệu năng của hai ngôn ngữ Python và Java. Ở đây chúng tôi chọn ngôn ngữ Python vì nó được dùng làm ngôn ngữ chính cho nhiều mô hình trí tuệ nhân tạo gần đây, và cũng thuận tiện cho chúng tôi khi so sánh hiệu năng cũng như đánh giá với các phương pháp học sâu khác da số được viết bằng Python.

Bảng \hyperref[tab:tab1]{2.1} bên dưới mô tả các kết quả thực nghiệm của chúng tôi với các độ đo \(H@K\) cùng với các kết quả thực nghiệm của các phương pháp khác được đề cập trong khảo  sát \cite{rossi2020knowledge}

\begin{table}[ht]
\caption{Kết quả phương pháp AnyBURL}
\label{tab:tab1}%
\begin{center}
\resizebox{\textwidth}{!}{%
\begin{tabular}{lllllllllllllllll}
\cline{2-17}
\multicolumn{1}{l|}{} &
  \multicolumn{4}{c|}{\textbf{FB15k}} &
  \multicolumn{4}{c|}{\textbf{FB15-237}} &
  \multicolumn{4}{c|}{\textbf{WN18}} &
  \multicolumn{4}{c|}{\textbf{WN18RR}} \\ \cline{2-17}
\multicolumn{1}{l|}{} &
  \multicolumn{1}{l|}{\textbf{H@1}} &
  \multicolumn{1}{l|}{\textbf{H@10}} &
  \multicolumn{1}{l|}{\textbf{MR}} &
  \multicolumn{1}{l|}{\textbf{MRR}} &
  \multicolumn{1}{l|}{\textbf{H@1}} &
  \multicolumn{1}{l|}{\textbf{H@10}} &
  \multicolumn{1}{l|}{\textbf{MR}} &
  \multicolumn{1}{l|}{\textbf{MRR}} &
  \multicolumn{1}{l|}{\textbf{H@1}} &
  \multicolumn{1}{l|}{\textbf{H@10}} &
  \multicolumn{1}{l|}{\textbf{MR}} &
  \multicolumn{1}{l|}{\textbf{MRR}} &
  \multicolumn{1}{l|}{\textbf{H@1}} &
  \multicolumn{1}{l|}{\textbf{H@10}} &
  \multicolumn{1}{l|}{\textbf{MR}} &
  \multicolumn{1}{l|}{\textbf{MRR}} \\ \hline
\multicolumn{1}{|l|}{\textbf{ComplEx}} &
  \multicolumn{1}{l|}{81.56} &
  \multicolumn{1}{l|}{90.53} &
  \multicolumn{1}{l|}{34} &
  \multicolumn{1}{l|}{0.848} &
  \multicolumn{1}{l|}{94.53} &
  \multicolumn{1}{l|}{95.50} &
  \multicolumn{1}{l|}{3623} &
  \multicolumn{1}{l|}{0.949} &
  \multicolumn{1}{l|}{25.72} &
  \multicolumn{1}{l|}{52.97} &
  \multicolumn{1}{l|}{202} &
  \multicolumn{1}{l|}{0.349} &
  \multicolumn{1}{l|}{42.55} &
  \multicolumn{1}{l|}{52.12} &
  \multicolumn{1}{l|}{4907} &
  \multicolumn{1}{l|}{0.458} \\ \hline
\multicolumn{1}{|l|}{\textbf{TuckER}} &
  \multicolumn{1}{l|}{72.89} &
  \multicolumn{1}{l|}{88.88} &
  \multicolumn{1}{l|}{39} &
  \multicolumn{1}{l|}{0.788} &
  \multicolumn{1}{l|}{94.64} &
  \multicolumn{1}{l|}{95.80} &
  \multicolumn{1}{l|}{510} &
  \multicolumn{1}{l|}{0.951} &
  \multicolumn{1}{l|}{25.90} &
  \multicolumn{1}{l|}{53.61} &
  \multicolumn{1}{l|}{162} &
  \multicolumn{1}{l|}{0.352} &
  \multicolumn{1}{l|}{42.95} &
  \multicolumn{1}{l|}{51.40} &
  \multicolumn{1}{l|}{6239} &
  \multicolumn{1}{l|}{0.459} \\ \hline
\multicolumn{1}{|l|}{\textbf{TransE}} &
  \multicolumn{1}{l|}{49.36} &
  \multicolumn{1}{l|}{84.73} &
  \multicolumn{1}{l|}{45} &
  \multicolumn{1}{l|}{0.628} &
  \multicolumn{1}{l|}{40.56} &
  \multicolumn{1}{l|}{94.87} &
  \multicolumn{1}{l|}{279} &
  \multicolumn{1}{l|}{0.646} &
  \multicolumn{1}{l|}{21.72} &
  \multicolumn{1}{l|}{49.65} &
  \multicolumn{1}{l|}{209} &
  \multicolumn{1}{l|}{0.31} &
  \multicolumn{1}{l|}{2.79} &
  \multicolumn{1}{l|}{49.65} &
  \multicolumn{1}{l|}{3936} &
  \multicolumn{1}{l|}{0.206} \\ \hline
\multicolumn{1}{|l|}{\textbf{RotatE}} &
  \multicolumn{1}{l|}{73.93} &
  \multicolumn{1}{l|}{88.10} &
  \multicolumn{1}{l|}{42} &
  \multicolumn{1}{l|}{0.791} &
  \multicolumn{1}{l|}{94.30} &
  \multicolumn{1}{l|}{96.02} &
  \multicolumn{1}{l|}{274} &
  \multicolumn{1}{l|}{0.949} &
  \multicolumn{1}{l|}{23.83} &
  \multicolumn{1}{l|}{53.06} &
  \multicolumn{1}{l|}{178} &
  \multicolumn{1}{l|}{0.336} &
  \multicolumn{1}{l|}{42.60} &
  \multicolumn{1}{l|}{57.35} &
  \multicolumn{1}{l|}{3318} &
  \multicolumn{1}{l|}{0.475} \\ \hline
\multicolumn{1}{|l|}{\textbf{ConvKB}} &
  \multicolumn{1}{l|}{11.44} &
  \multicolumn{1}{l|}{40.83} &
  \multicolumn{1}{l|}{324} &
  \multicolumn{1}{l|}{0.211} &
  \multicolumn{1}{l|}{94.89} &
  \multicolumn{1}{l|}{52.89} &
  \multicolumn{1}{l|}{202} &
  \multicolumn{1}{l|}{0.709} &
  \multicolumn{1}{l|}{13.98} &
  \multicolumn{1}{l|}{41.46} &
  \multicolumn{1}{l|}{309} &
  \multicolumn{1}{l|}{0.230} &
  \multicolumn{1}{l|}{5.63} &
  \multicolumn{1}{l|}{52.50} &
  \multicolumn{1}{l|}{3429} &
  \multicolumn{1}{l|}{0.249} \\ \hline
\multicolumn{1}{|l|}{\textbf{GAT}} &
  \multicolumn{1}{l|}{} &
  \multicolumn{1}{l|}{} &
  \multicolumn{1}{l|}{} &
  \multicolumn{1}{l|}{} &
  \multicolumn{1}{l|}{} &
  \multicolumn{1}{l|}{} &
  \multicolumn{1}{l|}{} &
  \multicolumn{1}{l|}{} &
  \multicolumn{1}{l|}{} &
  \multicolumn{1}{l|}{} &
  \multicolumn{1}{l|}{} &
  \multicolumn{1}{l|}{} &
  \multicolumn{1}{l|}{} &
  \multicolumn{1}{l|}{} &
  \multicolumn{1}{l|}{} &
  \multicolumn{1}{l|}{} \\ \hline
\multicolumn{1}{|l|}{\textbf{BURL}} &
  \multicolumn{1}{l|}{79.13} &
  \multicolumn{1}{l|}{82.30} &
  \multicolumn{1}{l|}{285} &
  \multicolumn{1}{l|}{0.824} &
  \multicolumn{1}{l|}{20.85} &
  \multicolumn{1}{l|}{42.40} &
  \multicolumn{1}{l|}{490} &
  \multicolumn{1}{l|}{0.311} &
  \multicolumn{1}{l|}{\textbf{93.86}} &
  \multicolumn{1}{l|}{\textbf{94.07}} &
  \multicolumn{1}{l|}{230} &
  \multicolumn{1}{l|}{0.955} &
  \multicolumn{1}{l|}{\textbf{44.22}} &
  \multicolumn{1}{l|}{54.23} &
  \multicolumn{1}{l|}{2528} &
  \multicolumn{1}{l|}{0.490} \\ \hline
Rank &
   &
   &
   &
   &
   &
   &
   &
   &
   &
   &
   &
   &
   &
   &
   &

\end{tabular}}
\end{center}
\end{table}

Bảng \hyperref[tab:tab2]{2.2} bên dưới mô tả các kết quả thực nghiệm của chúng tôi với hai chiến lược thêm tri thức mới vào đồ thị. Chúng tôi đánh giá trên hai phương diện là tổng số luật được sinh ra, và số luật có độ tin cậy \(>= 50\%\) và \(>= 80\%\).
\begin{table}[ht]
\caption{Kết quả phương pháp AnyBURL}
\label{tab:tab1}%
\begin{center}
\resizebox{\textwidth}{!}{%
\begin{tabular}{l|l|l|l|l|l|l|l|l|l|l|l|l|}
\cline{2-13}
                                 & \multicolumn{3}{l|}{FB15k}                & \multicolumn{3}{l|}{FB15k-237}        & \multicolumn{3}{l|}{WN18}                 & \multicolumn{3}{l|}{WN18RR}            \\ \cline{2-13} 
                                 & num\_rule & 50\%          & 80\%          & num\_rule & 50\%         & 80\%       & num\_rule & 50\%           & 80\%         & num\_rule & 50\%         & 80\%        \\ \hline
\multicolumn{1}{|l|}{onl-to-off} & 1011      & 416(41,14\%)  & 284(28, 09\%) & 1120      & 244(21,79\%) & 95(8,48\%) & 533       & 270(38, 46 \%) & 240(34,19\%) & 439       & 110(25,05\%) & 83(18,91)   \\ \hline
\multicolumn{1}{|l|}{onl-to-onl} & 1367      & 1185(86,69\%) & 481(35,18\%)  & 756       & 660(87,30\%) & 162(21,43) & 260       & 252(96,92\%)   & 225(86,54\%) & 106       & 102(96,22)   & 85(81,19\%) \\ \hline
\end{tabular}}
\end{center}
\end{table}
\section{Kết luận}
Trong phần này chúng tôi đưa ra các kết quả đạt được của mô hình chúng tôi, chúng tôi cố gắng tìm hiểu các đặc trưng của các bộ dữ liệu tương ứng để cố gắng lý giải thích tại sao mô hình của chúng tôi hoặc các công trình khác có được kết quả tốt trên tập dữ liệu tương ứng đó. Những kết quả của hai đề xuất của chúng tôi cũng như các dịnh hướng nghiên cứu của chúng tôi trong tương lai.

Mặc dù kết quả chúng tôi cho thấy phương pháp của chúng tôi có hiệu suất tương đương với các mô hình học sâu hiện đại(state-of-art) và có ưu thế vượt trội trong thời gian đào tạo khoảng 17 phút so với thời gian hàng giờ của phương pháp học sâu khác nhưng không phải là các mô hình deep learning này không đáng nghiên cứu. Chúng tôi cũng nhận thấy đối với những tập dữ liệu khó như FB15-237 hay WN18RR phương pháp của chúng tôi thường cho kết quả không tốt do các quan hệ tương tự hay nghịch đảo không không xuất hiện trong ví dụ đào tạo nên chúng tôi khó tạo ra các luật đủ tốt để có thể khái quát hóa trên toàn bộ đồ thị đẫn đến các kết quả không tốt. Ngược lại đối với các phương pháp dựa trên học sâu lại có ưu thế rất lớn trong các tập dữ liệu này do có thể dễ đàng tính toán độ gần của các luật mới cần đánh giá so với các luật đã học từ đó có một kết quả khá tốt. Do đó chúng tôi cũng sẽ tiếp tục nghiên cứu các phương pháp học sâu và sẽ dùng phương pháp này làm đường cơ sở(base line) để so sánh với các nghiên cứu của chúng tôi trong tương lai. Một điểm yếu nữa của mô hình đựa trên luật của chúng tôi là mặc dù thời gian học là vượt trội nhưng thời gian để tính toán đưa ra đự đoán khá lâu do phải duyệt qua tất cả các luật được sinh ra mới có thể đưa ra dự đoán. Không giống như các phương pháp nhúng đồ thị khác thao tác này có thể dễ dàng tính toán.

Đối với hai thuật toán mở rộng của chúng tôi trong việc thêm tri thức mới vào đồ thị chúng tôi nhận thấy rằng là vượt trội hoàn toàn so với các phương pháp học sâu. Ở các phương pháp học sâu điều này đường như chưa được ai chú trọng nghiên cứu mặc dù thời gian đào tạo một mô hình là tương đối mất thời gian. Khi có tri thức mới hầu hết các mô hình phải đào tạo lại toàn bộ điều này khá lãng phí. Chúng tôi cũng xem đây là mục tiêu tiếp theo cho chúng tôi khi nghiên cứu các mô hình học sâu trong tương lai. Gần đây nhánh học tăng cường(Reinforcement learning) khá phát triển và nhóm tác giả Meilicke, Christian and Chekol \cite{meilicke2020reinforced} gần đây cũng đã có 1 nghiên cứu để tối ưu hóa lại phương pháp AnyBURL này. Chúng tôi cũng có ý định nghiên cứu về hướng này và cố gắng báo cáo lại trong một tương lai gần. 

\section{Bố cục của báo cáo}

Nội dung của báo cáo tối thiểu 50 trang khổ A4 và không nên vượt quá 100 trang (không kể các trang bìa, lời cám ơn, mục lục, tài liệu tham khảo \ldots) theo trình tự như sau:

\begin{itemize}
\item MỞ ĐẦU (thường đặt tên là ``Giới thiệu''): Trình bày lý do chọn đề tài, mục đích, đối tượng và phạm vi nghiên cứu.
Mô tả bài toán mà đề tài giải quyết.
Bài toán này có gì hay?
Tại sao lại cần giải quyết bài toán này?
Bài toán này có gì khó?
Có những hướng nào để giải quyết bài toán này?
Những hướng giải quyết trước đây có những vấn đề gì chưa giải quyết được?
Các câu hỏi nghiên cứu mà đề tài trả lời hoặc những vấn đề mà đề tài sẽ giải quyết.
Các đóng góp của đề tài.

\item TỔNG QUAN (thường đặt tên là ``Các công trình liên quan''): Phân tích đánh giá các hướng nghiên cứu đã có của các tác giả trong và ngoài nước liên quan đến đề tài; nêu những vấn đề còn tồn tại (những vấn đề nào mà các công trình khác chưa giải quyết được); chỉ ra những vấn đề mà đề tài cần tập trung, nghiên cứu giải quyết.

\item NGHIÊN CỨU THỰC NGHIỆM HOẶC LÝ THUYẾT (thường đặt tên là ``Phương pháp đề xuất''): Trình bày cơ sở lý thuyết, lý luận, giả thiết khoa học và phương pháp nghiên cứu đã được sử dụng trong đề tài.

Nếu đề xuất hướng giải quyết mới, mô hình mới thì cần mô tả chi tiết cách giải quyết của mình (chi tiết tới mức người khác có thể dựa vào phần này mà cài đặt lại được đúng hoàn toàn phương pháp của mình đề ra).

\item TRÌNH BÀY, ĐÁNH GIÁ BÀN LUẬN VỀ CÁC KẾT QUẢ (thường đặt tên là ``Kết quả thí nghiệm''): Mô tả các kết quả nghiên cứu khoa học hoặc kết quả thực nghiệm.
Đối với  đề tài ứng dụng có kết quả là sản phẩm phần mềm phải có hồ sơ thiết kế, cài đặt,\ldots theo một trong các mô hình đã học (UML,\ldots).

Thông thường cần mô tả môi trường thí nghiệm trước như sử dụng dữ liệu nào, dùng độ đo nào để đánh giá, môi trường chạy thí nghiệm (cấu hình máy nếu cần phân tích thông tin về thời gian chạy thực nghiệm). Sau đó, nêu kết quả thực nghiệm, bàn luận và giải thích kết quả.

\item KẾT LUẬN VÀ HƯỚNG PHÁT TRIỂN (thường đặt tên là ``Kết luận''): Trình bày những kết quả đạt được, những đóng góp mới và những đề xuất mới, kiến nghị về những hướng nghiên cứu tiếp theo.

\item DANH MỤC TÀI LIỆU THAM KHẢO: Chỉ bao gồm các tài liệu được trích dẫn, sử dụng và đề cập tới để bàn luận trong báo cáo.
Phần này các bạn chuẩn bị 1 file BIB để lưu các tài liệu trích dẫn.
Khi các bạn trích dẫn một tài liệu nào đó, LaTeX sẽ tự động thêm vào danh mục tài liệu tham khảo giúp các bạn.
Các bạn xem hướng dẫn cách trích dẫn ở chương sau.

\item PHỤ LỤC: Phần này bao gồm nội dung cần thiết nhằm minh họa hoặc hỗ trợ cho nội dung báo cáo như số liệu, mẫu biểu, tranh ảnh,\ldots Phụ lục không được dày hơn phần chính của báo cáo.
Nếu có công trình công bố thì để vào phần phụ lục này.
\end{itemize}

\section{Bảng biểu, hình vẽ, phương trình}

%Những qui định dưới này các bạn có thể bỏ qua hoặc đọc để hiểu thêm.
%Những định dạng này LaTeX đều tự động giúp các bạn.
%Các bạn xem hướng dẫn chi tiết hơn ở chương sau.

Việc đánh số bảng biểu, hình vẽ, phương trình phải gắn với số chương; ví dụ hình 3.4 có nghĩa là hình thứ 4 trong Chương 3.
Mọi đồ thị, bảng biểu, hình vẽ lấy từ các nguồn khác phải được trích dẫn đầy đủ.

\subsection{Bảng biểu, hình vẽ}



Đầu đề của bảng biểu ghi phía trên bảng, đầu đề của hình vẽ ghi phía dưới hình.

Thông thường, những bảng ngắn và đồ thị phải đi liền với phần nội dung đề cập tới các bảng và đồ thị này ở lần thứ nhất.
Các bảng dài có thể để ở những trang riêng nhưng cũng phải tiếp theo ngay phần nội dung đề cập tới bảng này ở lần đầu tiên.
Các bảng rộng vẫn nên trình bày theo chiều đứng dài 297mm của trang giấy, chiều rộng của trang giấy có thể hơn 210mm.
Chú ý gấp trang giấy sao cho số và đầu đề của hình vẽ hoặc bảng vẫn có thể nhìn thấy ngay mà không cần mở rộng tờ giấy.
Tuy nhiên hạn chế sử dụng các bảng quá rộng này.

Đối với những trang giấy có chiều đứng hơn 297mm (bản đồ, bản vẽ,\ldots) thì có thể để trong một phong bì cứng đính bên trong bìa sau của báo cáo.

Các hình vẽ phải sạch sẽ bằng mực đen để có thể sao chụp lại; có đánh số và ghi đầy đủ đầu đề, cỡ chữ phải bằng cỡ chữ sử dụng trong báo cáo.

Khi đề cập đến các bảng biểu và hình vẽ phải nêu rõ số của hình và bảng biểu đó, ví dụ ``... được nêu trong Bảng 4.1'' hoặc ``xem Hình 3.2'' mà không được viết ``… được nêu trong bảng dưới đây'' hoặc ``trong đồ thị của X và Y sau''.

\subsection{Phương trình toán học}

Việc trình bày phương trình toán học trên một dòng đơn hoặc dòng kép tùy ý, tuy nhiên phải thống nhất trong toàn báo cáo.

Khi ký hiệu xuất hiện lần đầu tiên thì phải giải thích và đơn vị tính phải đi kèm ngay trong phương trình có ký hiệu đó.
Nếu cần thiết, danh mục của tất cả các ký hiệu, chữ viết tắt và nghĩa của chúng cần được liệt kê và để ở phần đầu của báo cáo.

Tất cả các phương trình cần được đánh số và để trong ngoặc đơn đặt bên phía lề phải.
Nếu một nhóm phương trình mang cùng một số thì những số này cũng được để trong ngoặc, hoặc mỗi phương trình trong nhóm phương trình (5.1) có thể được đánh số là (5.1.1), (5.1.2), (5.1.3).

\section{Viết tắt}

\textbf{Không lạm dụng việc viết tắt} trong báo cáo.
Chỉ viết tắt những từ, cụm từ hoặc thuật ngữ được sử dụng nhiều lần trong báo cáo.
Không viết tắt những cụm từ  dài, những mệnh đề; không viết tắt những cụm từ ít xuất hiện trong báo cáo.
Nếu cần viết tắt những từ thuật ngữ, tên các cơ quan, tổ chức,\ldots thì được viết tắt sau lần viết thứ nhất có kèm theo chữ viết tắt trong ngoặc đơn.
Nếu báo cáo có nhiều chữ viết tắt thì phải có bảng danh mục các chữ viết tắt (xếp theo thứ tự ABC) ở phần đầu báo cáo.

Nhắc lại: \textbf{không lạm dụng việc viết tắt} trong báo cáo.
Khi các bạn sử dụng từ viết tắt, người đọc sẽ phải lật lại những phần đã đọc, để tìm lại xem từ viết tắt đó nghĩa là gì.
Việc này sẽ làm chậm tốc độ đọc và sẽ khiến người đọc khó theo dõi báo cáo của bạn hơn.
Nếu có thể, hạn chế hoàn toàn việc dùng viết tắt.

\section{Tài liệu tham khảo và cách trích dẫn}

Mọi ý kiến, khái niệm có ý nghĩa, mang tính chất gợi ý không phải của riêng tác giả và mọi tham khảo khác phải được trích dẫn và chỉ ra nguồn trong danh mục Tài liệu tham khảo của báo cáo. Nguồn được trích dẫn phải được liệt kê chính xác trong danh mục Tài liệu tham khảo.

Việc trích dẫn, tham khảo chủ yếu nhằm thừa nhận nguồn của những ý tưởng có giá trị giúp người đọc theo được mạch suy nghĩ của tác giả, không làm trở ngại việc đọc.


Không trích dẫn những kiến thức phổ biến, mọi người đều biết cũng như không làm báo cáo nặng nề với những tham khảo trích dẫn.

Nếu không có điều kiện tiếp cận được một tài liệu gốc mà phải trích dẫn thông qua một tài liệu khác thì phải nêu ra trích dẫn này, đồng thời tài liệu gốc đó không được liệt kê trong danh mục tài liệu tham khảo của báo cáo.

% Phương pháp đề xuất
\chapter{Hướng dẫn sử dụng template}
\label{Chapter3}

\section{Trích dẫn tài liệu}

Dùng lệnh $\textbackslash cite$ để trích dẫn một hoặc nhiều tài liệu tham khảo.
Tài liệu tham khảo có thể là trang web~\cite{Listings,HDLVThS}, bài báo khoa học~\cite{1994-Cavnar}, sách~\cite{1984-TeX-Knuth,2006-DDien,2006-NPTV}, bài tạp chí~\cite{1989-TED} hoặc các nguồn tham khảo khác. 
Lưu ý khi trích dẫn tài liệu tham khảo, cần viết câu sao cho bỏ phần trong cặp ngoặc vuông đi thì câu vẫn đầy đủ ý nghĩa.
Ví dụ, thay vì viết ``Nghiên cứu~\cite{1989-TED} chỉ ra rằng ... '' thì nên viết ``Nghiên cứu của Zhang~\cite{1989-TED} chỉ ra rằng ...''.
Một ví dụ khác, thay vì viết ``... như trong công trình nghiên cứu~\cite{1994-Cavnar}.'' thì nên viết ``... như trong công trình nghiên cứu của Cavnar và Trenkle~\cite{1994-Cavnar}.''

\section{Chèn mã nguồn}

Để chèn mã nguồn, cần dùng package listings~\cite{Listings}:

\begin{lstlisting}
\usepackage{listings}
\end{lstlisting}

Mã nguồn có thể được chèn trực tiếp như sau:

\begin{lstlisting}[language=Python]
print "Hello, World!"
\end{lstlisting}

hoặc chèn thông qua tập tin chứa mã nguồn trong thư mục $SourceCode$ như sau:

\lstinputlisting[language=C++]{SourceCode/hello.cpp}

Để chèn mã giả, cần dùng package algorithm~\cite{Algorithm}:

\begin{lstlisting}
\usepackage{algorithm}
\end{lstlisting}

Có thể chèn mã giả vào như sau:

\begin{algorithm}
\caption{My algorithm}\label{euclid}
\begin{algorithmic}[1]
\Procedure{MyProcedure}{}
\State $\textit{stringlen} \gets \text{length of }\textit{string}$
\State $i \gets \textit{patlen}$
\BState \emph{top}:
\If {$i > \textit{stringlen}$} \Return false
\EndIf
\State $j \gets \textit{patlen}$
\BState \emph{loop}:
\If {$\textit{string}(i) = \textit{path}(j)$}
\State $j \gets j-1$.
\State $i \gets i-1$.
\State \textbf{goto} \emph{loop}.
\State \textbf{close};
\EndIf
\State $i \gets i+\max(\textit{delta}_1(\textit{string}(i)),\textit{delta}_2(j))$.
\State \textbf{goto} \emph{top}.
\EndProcedure
\end{algorithmic}
\end{algorithm}

\section{Hình ảnh}

Để chèn hình ảnh, cần dùng package graphicx~\cite{Figures}:

\begin{lstlisting}
\usepackage{graphicx}
\end{lstlisting}

Hình \ref{fig:vd1}, hình \ref{fig:vd2} là một số ví dụ về chèn hình ảnh.

\begin{figure}[htp]
\centering
\includegraphics[width=6 cm]{images/logo-khtn.png}
\caption{Hình ví dụ 1}
\label{fig:vd1}
\end{figure}

\begin{figure*}[htp]
\centering
\includegraphics[width=40 mm]{images/logo-khtn.png}
\caption{Hình ví dụ 2}
\label{fig:vd2}
\end{figure*}

\section{Bảng biểu}

Để tạo bảng biểu, tham khảo thêm tại sharelatex.com~\cite{Tables}.
Bảng~\ref{tab:vd1} là một ví dụ về bảng.
Ngoài ra, có một số tool online~\footnote{Như trang http://www.tablesgenerator.com/} có thể được dùng để tạo bảng biểu một cách trực quan.

\begin{table}[ht]
\caption{Bảng ví dụ 1}

\label{tab:vd1}%
\begin{center}
\begin{tabular}{ |p{3cm}||p{3cm}|p{3cm}|p{3cm}|  }
 \hline
 \multicolumn{4}{|c|}{Country List} \\
 \hline
 Country Name     or Area Name& ISO ALPHA 2 Code &ISO ALPHA 3 Code&ISO numeric Code\\
 \hline
 Afghanistan   & AF    &AFG&   004\\
 Aland Islands&   AX  & ALA   &248\\
 Albania &AL & ALB&  008\\
 Algeria    &DZ & DZA&  012\\
 American Samoa&   AS  & ASM&016\\
 Andorra& AD  & AND   &020\\
 Angola& AO  & AGO&024\\
 \hline
\end{tabular}
\end{center}
\end{table}

\section{Công thức}

Công thức có thể chèn vào trong cùng một dòng như $ \sqrt{a^2+b^2} $ hoặc nằm trên dòng riêng như công thức~\ref{eu_eqn}.

\begin{equation} \label{eu_eqn}
x = a_0 + \frac{1}{a_1 + \frac{1}{a_2 + \frac{1}{a_3 + a_4}}}
\end{equation}

% Kết quả thí nghiệm
\chapter{Kết quả thí nghiệm}
\label{Chapter4}

Trong phần này chúng tôi mô tả lại các bộ dữ liệu mà chúng tôi đùng để thực nghiệm đánh giá phương pháp của chúng tôi cùng với so sánh với các kết quả khác của các công trình nổi bật khác được báo cáo trong \hyperref[tab:tab1]{bảng 2.1}. Ngoài ra chúng tôi cũng cố gắng đánh giá hai đề xuất của chúng tôi trong việc thêm một lượng tri thức mới vào đồ thị bằng cách chúng tôi xem tập test là một lượng tri thức mới cần thêm vào và dùng tập vadidation để đánh giá lại phương pháp của chúng tôi. Kết quả chi tiết được báo cáo ở \hyperref[tab:tab1]{bảng 2.2}

\begin{table}[htbp]
	\begin{center}
		\resizebox{\textwidth}{!}{%
			\begin{tabular}{lllllll}
				\hline
				&          &           & \multicolumn{4}{l}{\# Edges}    \\ \cline{4-7}

				Dataset   & Entities & Relations & Training & Validation & Test   & Total   \\ \hline
				FB15k     & 14,951   & 1,345     & 483,142  & 50,000     & 59,071 & 592,213 \\
				FB15k-237 & 14,541   & 237       & 272,115  & 17,535     & 20,466 & 310,116 \\
				WN18      & 40,943   & 18        & 141,442  & 5,000      & 5,000  & 151,442 \\
				WN18RR    & 40,943   & 11        & 86,835   & 3034       & 3134   & 93,003 \\
				\hline
			\end{tabular}
		}
		\caption{Thông tin các tập dữ liệu}
		\label{tab:dataset}
	\end{center}
\end{table}


\section{Các tập dữ liệu huấn luyện} \label{datasets}
Trong thí nghiệm của chúng tôi, chúng tôi thực hiện trên bốn tập dữ liệu phổ biến là FB15k, FB15-237, WN18 và WN18RR.Các bộ dữ liệu này là tập hợp các bộ ba (triple) \(\langle head, relation, tail \rangle\) biểu thị cho thực thể đầu có một mối quan hệ với thực thể cuối.

\subsection{Bộ dữ liệu FB15k}
Bộ dữ liệu này được tạo bởi nhóm nghiên cứu A. Bordes, N. Usunier, A. Garcia-Duran, J. Weston, and O. Yakhnenko \cite{bordes2013translating} bằng cách trích xuất từ bộ dữ liệu Wikilinks database \footnote{https://code.google.com/archive/p/wiki-links/}.Wikilinks database thu thập các siêu liên kết(hyperlinks) đến Wikipedia gồm 40 triệu lượt đề cập trên 3 triệu thực thể, họ trích xuất tất cả các dữ kiện liên quan đến một thực thể nhất định có hơn 100 lần được đề cập đến bởi các tài liệu khác cùng với tất cả các dữ kiện liên quan đến thực thể đó (bao gồm cả những thực thể con được nhắc đến trong tài liệu Wikipedia đó), ngoại trừ những thonng tin như: ngày tháng, danh từ riêng, v.v ... Họ cũng chuyển đổi các đỉnh có bậc \(n\) được biểu diễn thành các nhóm các cạnh nhị phân tức là liệt kê các cạnh và quan hệ của mọi đỉnh. Tập dữ liệu được chia ngẫu nhiên thành 3 tập: tập training với 1345 relations, 14834 head entities và 14903 tail entities, tập test gồm 916 relations, 11886 head entities, và 11285 tail entities, tập validation gồm 961 relations, 12297 head entities, và 11825 tail entities.

\subsection{Bộ dữ liệu FB15k-237}
Bộ dữ liệu này là một tập hợp con của FB15k được xây dựng bởi Toutanova và Chen \cite{toutanova2015observed} lấy cảm hứng từ quan sát rằng FB15k bao gồm dữ liệu thử nghiệm được các mô hình nhìn thấy tại thời điểm đào tạo (test leekage). Trong FB15k, vấn đề này là do sự hiện diện của các quan hệ gần giống nhau hoặc nghịch đảo của nhau.FB15k-237 được xây dựng để trở thành một tập dữ liệu thách thức hơn: các tác giả đã chọn các dữ kiện liên quan đến 401 quan hệ xuất hiện nhiều nhất và loại bỏ tất cả các quan hệ tương đương hoặc nghịch đảo. Họ cũng đảm bảo rằng không có thực thể nào được kết nối trong tập huấn luyện cũng được liên kết trực tiếp trong tập test và validation.Tập training gồm 237 relations, 13781 head entities, và 13379 tail entities, tập test gồm 223 relations, 7652 head entities, và 5804 tail entities, tập vadition gồm 224 relations, 8171 head entities, and 6376 tail entities.

\subsection{Bộ dữ liệu WN18}
Bộ dữ liệu này được giới thiệu bởi các tác giả của TransE \cite{bordes2013translating}, được trích xuất từ WordNet\footnote{https://wordnet.princeton.edu/}, một bản thể học ngôn ngữ KG có nghĩa là cung cấp một từ điển/từ đồng nghĩa để hỗ trợ NLP và phân tích văn bản tự động. Trong WordNet, các thực thể tương ứng với các tập hợp (\textit{word senses}) và các quan hệ đại diện cho các kết nối từ vựng của chúng (ví dụ: “hypernym”). Để xây dựng WN18, các tác giả đã sử dụng WordNet làm điểm bắt đầu và sau đó lặp đi lặp lại lọc ra các thực thể và mối quan hệ với quá ít lần được đề cập. Tập dữ liệu được chia ngẫu nhiên thành 3 tập: tập training với 18 relations, 40504 head entities, và 40551 tail entities, tập test gồm 18 relations, 4262 head entities, and 4338 tail entities, tập vadiation gồm 18 relations, 4349 head entities, and 4263 tail entities.

\subsection{Bộ dữ liệu WN18RR}
Bộ dữ liệu này là một tập hợp con của WN18 được xây dựng bởi DeŠmers et al.\cite{dettmers2017convolutional}, cũng là người giải quyết vấn đề rò rỉ thử nghiệm (test leakage)trong WN18. Để giải quyết vấn đề đó, họ xây dựng tập dữ liệu WN18RR thách thức hơn nhiều bằng cách áp dụng một phương pháp tương tự được sử dụng cho FB15k-237 \cite{toutanova2015observed}. Training gồm 11 relations, 39610 head entities, và 31881 tail entities, tập test gồm 11 relations, 2958 head entities, và 2619 tail entities, tập vadition gồm 11 relations, 2851 head entities, and 2575 tail entities.

\section{Kết quả thực nghiệm}
Trong phần này chúng tôi mô tả lại các phương pháp đánh giá(độ do), môi trường thực hiện cũng như các tập dữ liệu mà chúng tôi sử dụng để dánh giá phương pháp của mình. Các phương pháp đánh giá(độ do) này cũng phổ biến nó được đánh giá cho hầu hết các mô hình dự đoán liên kết trên đồ thị. Chúng tôi tiến hành so sánh với bốn phương pháp nổi bật khác được báo cáo trong \cite{rossi2020knowledge}.
\subsection{Các độ đo}
\textit{Mean Rank} (MR). Đây là giá trị trung bình của rank thu được cho một dự đoán chính xác. Càng nhỏ thì mô hình càng chính xác:
\[MR = \frac{1}{\mid Q \mid} \sum_{q ~\in~ Q} rank(q) \]
Trong đó \(\mid Q \mid\) là độ lớn của tập hợp các câu hỏi bằng độ lớn của tập test hoặc vadidation. Khi dự đoán chúng tôi dự đoán cả head và tail cho một dòng tương ứng trong tập dữ liệu thử nghiệm. Ví dụ chún tôi sẽ dự đoán \(\langle ?,~ relation,~ tail \rangle\) và \(\langle head,~ relation,~ ?\rangle\) cho 1 dòng tương ứng, \(q\) thể hiện cho câu hỏi chúng tôi dự đoán và \(rank(q)\) thể hiện cho kết quả đúng của câu hỏi đứng ở vị trí thứ mấy trong xếp hạng của chúng tôi sau đó lấy trung bình rank của các dự đoán head và tail. Rõ ràng độ đo này nằm giữa \([1, \mid \text{số lượng các entity} \mid]\) do có tối da \(n\) cạnh nối 1 đỉnh tới \(n-1\) đỉnh còn lại và thêm cạnh nối tới chính đỉnh nó(cạnh khuyên). Và độ đo này đễ bị ảnh hưởng bởi nhiễu vì có những quan hệ có những thực thể được xếp hạng gần cuối. Để giải quyết vấn đề này nhóm chúng tôi và các nhóm nghiên cứu khác sử dụng thêm độ đo Mean Reciprocal Rank (MMR).

\textit{Mean Reciprocal Rank} (MMR). Đây là xếp hạng đối ứng trung bình, là nghịch đảo của giá trị trung bình của rank thu được cho một dự đoán chính xác ở trên. Và càng lớn thì mô hình càng chính xác. Do độ đo này lấy nghịch đảo của các rank nên tránh dược vấn đề nhiễu của độ đo MR ở trên.
\[MRR =\frac{1}{\mid Q \mid} \sum_{q~ \in ~Q} \frac{1}{rank(q)}\]

\textit{Hit@K} (H@K). Đó là tỷ lệ các dự đoán đúng mà rank nhỏ hơn hoặc bằng ngưỡng \(K\):
\[H@K = \frac{\mid {q ~\in ~Q~: rank(q) \leq K} \mid}{\mid Q \mid}\]
\subsection{Kết quả}
Như đã nói trước đây với mô hình dựa trên luật của chúng tôi hoàn toàn có thể thực hiện trên một laptop với cấu hình thông thường. Trong thí nghiệm của chúng tôi cấu hình máy để thực thi như sau: T480, core i5 8th Gen, ram 16Gb, 4 core 8 thread. Mã nguồn thực thi được viết bằng ngôn ngữ Python phiên bản 3.6 và dùng các hàm hỗ trợ có sẵn trong Python với không một thư viện bên thứ ba nào. Thí nghiệm được thực hiện với bốn tập dữ liệu phổ biến là FB15k, FB15-237, WN18 và WN18RR. Thông tin chi tiết các bộ dữ liệu này được mô tả ở phần \hyperref[datasets]{các tập dữ liệu huấn luyện}.

Như mô tả ở phần \hyperref[algorithm2]{thuật toán AnyBURL} thuật toán này sẽ học các luật được sinh ra trong một khoảng thời gian nhất định do người dùng cấu hình. Ở đây chúng tôi chọn cấu hình thời gian là 1000 giây tương đương khoảng 17 phút đào tạo, với độ bão hòa(SAT) \(0.85\), độ tin cậy Q \(0.05\), kích thước mẫu S (\(\frac{1}{10}~ \text{tập huấn luyện}\)). Với cấu hình như vậy mô hình phiên bản Python của chúng tối cho kết quả tương đương với phiên bản Java nhóm tác giả Meilicke, Christian et al. \cite{burl} với cấu hình tương tự nhưng thời gian training là 100 giây. Sự khác biệt về thời gian học tập ở đây chủ yếu là do hiệu năng của hai ngôn ngữ Python và Java. Ở đây chúng tôi chọn ngôn ngữ Python vì nó được dùng làm ngôn ngữ chính cho nhiều mô hình trí tuệ nhân tạo gần đây, và cũng thuận tiện cho chúng tôi khi so sánh hiệu năng cũng như đánh giá với các phương pháp học sâu khác da số được viết bằng Python.

\hyperref[tab:tab1]{Bảng 4.1}, và \hyperref[tab:tab2]{bảng 4.2} bên dưới mô tả các kết quả thực nghiệm của chúng tôi với các độ đo \(H@K\) cùng với các kết quả thực nghiệm của các phương pháp khác được đề cập trong khảo  sát \cite{rossi2020knowledge}

\begin{table}[ht]
\begin{center}
\resizebox{\textwidth}{!}{%
\begin{tabular}{l|l|l|l|l|l|l|l|l|}
\cline{2-9}
                                       & \multicolumn{4}{c|}{\textbf{FB15k}}                   & \multicolumn{4}{c|}{\textbf{FB15k-237}}                   \\ \cline{2-9} 
                                       & \textbf{H@1} & \textbf{H@10} & \textbf{MR} & \textbf{MRR} & \textbf{H@1} & \textbf{H@10} & \textbf{MR} & \textbf{MRR} \\ \hline
\multicolumn{1}{|l|}{ComplEx} & 81.56        & 90.53         & 34          & 0.848        & 25.72        & 52.97         & 202        & 0.349        \\ \hline
\multicolumn{1}{|l|}{TuckER}  & 72.89        & 88.88         & 39          & 0.788        & 25.90        & 53.61         & 162         & 0.352        \\ \hline
\multicolumn{1}{|l|}{TransE}  & 49.36        & 84.73         & 45          & 0.628        & 21.72        & 49.65         & 209         & 0.31        \\ \hline
\multicolumn{1}{|l|}{RoteE}   & 73.93        & 88.10         & 42          & 0.791        & 23.83        & 53.06         & 178         & 0.336        \\ \hline
\multicolumn{1}{|l|}{ConvKB}  & 59.46        & 84.94         & 51         & 0.688        & 21.90        & 47.62         & 281         &0.305        \\ \hline
\multicolumn{1}{|l|}{\textbf{KBGAT}}     &              &               &             &              & 36.06     &    58.32   &  211  &    0.4353  \\ \hline
\multicolumn{1}{|l|}{\textbf{AnyBURL}}    & \textbf{79.13}        & \textbf{82.30}         & \textbf{285}         & \textbf{0.824}        & \textbf{20.85}        & \textbf{42.40}         & \textbf{490}         & \textbf{0.311}        \\ \hline
\end{tabular}
}
\caption{Kết quả phương pháp đề xuất trên FB15k, FB15k-237}
\label{tab:tab1}%
\end{center}
\end{table}

\begin{table}[ht]
\begin{center}
\resizebox{\textwidth}{!}{%
\begin{tabular}{l|l|l|l|l|l|l|l|l|}
\cline{2-9}
                                       & \multicolumn{4}{c|}{\textbf{WN18}}                              & \multicolumn{4}{c|}{\textbf{WN18RR}}                            \\ \cline{2-9} 
                                       & \textbf{H@1}   & \textbf{H@10}  & \textbf{MR}  & \textbf{MRR}   & \textbf{H@1}   & \textbf{H@10}  & \textbf{MR}  & \textbf{MRR}   \\ \hline
\multicolumn{1}{|l|}{ComplEx} & 94.53          & 95.50          & 3623         & 0.349          & 42.55          & 52.12          & 4909         & 0.458          \\ \hline
\multicolumn{1}{|l|}{TuckER}  & 94.64          & 95.80          & 510          & 0.951          & 42.95          & 51.40          & 6239         & 0.459          \\ \hline
\multicolumn{1}{|l|}{TransE}  & 40.56          & 94.87          & 279          & 0.646          & 2.79           & 94.87          & 279          & 0.646          \\ \hline
\multicolumn{1}{|l|}{RoteE}   & 94.30          & 96.02          & 274          & 0.949          & 42.60          & 57.35          & 3318         & 0.475          \\ \hline
\multicolumn{1}{|l|}{ConvKB}  & 93.89          & 95.68          & 413          & 0.945          & 38.99          & 50.75          & 4944         & 0.427          \\ \hline
\multicolumn{1}{|l|}{\textbf{KBGAT}}     &                &        &        &                &       25.74         &        52.34         &      2773       &  35.24           \\ \hline
\multicolumn{1}{|l|}{\textbf{AnyBURL}}    & \textbf{93.96} & \textbf{95.07} & \textbf{230} & \textbf{0.955} & \textbf{44.22} & \textbf{54.40} & \textbf{490} & \textbf{0.490} \\ \hline
\end{tabular}
}
\caption{Kết quả phương pháp AnyBURL trên WN18, WN18RR}
\label{tab:tab2}%
\end{center}
\end{table}

Bảng \hyperref[tab:tab3]{4.3} bên dưới mô tả các kết quả thực nghiệm của chúng tôi với hai chiến lược thêm tri thức mới vào đồ thị. Chúng tôi đánh giá trên tổng số luật được sinh ra, và số luật có độ tin cậy \(>= 50\%\) và \(>= 80\%\).
\begin{table}[ht]
\begin{center}
\begin{tabular}{ll|l|l|}
\cline{3-4}
                                                          &                 & \textbf{online-to-offline} & \textbf{online-to-online} \\ \hline
\multicolumn{1}{|c|}{\multirow{3}{*}{\textbf{FB15k}}}     & num rule        & 1011                       & 1367                      \\ \cline{2-4} 
\multicolumn{1}{|c|}{}                                    & confidence 50\% & 416 (41,14\%)              & 1185 (86,69\%)            \\ \cline{2-4} 
\multicolumn{1}{|c|}{}                                    & confidence 80\% & 284 (28, 09\%)             & 481 (35,18\%)             \\ \hline
\multicolumn{1}{|l|}{\multirow{3}{*}{\textbf{FB15k-237}}} & num rule        & 1120                       & 756                       \\ \cline{2-4} 
\multicolumn{1}{|l|}{}                                    & confidence 50\% & 244 (21,79\%)              & 660 (87,30\%)             \\ \cline{2-4} 
\multicolumn{1}{|l|}{}                                    & confidence 80\% & 95 (8,48\%)                & 162 (21,43\%)             \\ \hline
\multicolumn{1}{|l|}{\multirow{3}{*}{\textbf{WN18}}}      & num rule        & 533                        & 260                       \\ \cline{2-4} 
\multicolumn{1}{|l|}{}                                    & confidence 50\% & 270 (38, 46 \%)            & 252 (96,92\%)             \\ \cline{2-4} 
\multicolumn{1}{|l|}{}                                    & confidence 80\% & 240 (34,19\%)              & 225 (86,54\%)             \\ \hline
\multicolumn{1}{|l|}{\multirow{3}{*}{\textbf{WN18RR}}}    & num rule        & 439                        & 106                       \\ \cline{2-4} 
\multicolumn{1}{|l|}{}                                    & confidence 50\% & 110 (25,05\%)              & 102 (96,22\%)             \\ \cline{2-4} 
\multicolumn{1}{|l|}{}                                    & confidence 80\% & 83 (18,91\%)               & 85 (81,19\%)              \\ \hline
\end{tabular}
\caption{Kết quả hai chiến lược thêm tri thức mới}
\label{tab:tab3}%
\end{center}
\end{table}

% Kêt luận
\chapter{Kết luận}
\label{Chapter5}

Trong phần này chúng tôi sẽ trình bày các kết quả đạt được của mô hình chúng tôi, cũng như những phân tích của chúng tôi
trên các kết quả của các tập dữ liệu khác nhau để giải thích những điểm tốt và điểm cần cải thiện trên mô hình của chúng tôi trên tập dữ liệu đó. Từ đó chúng tôi xác định những hướng nghiên cứu để cải tiến trong tương lai.
% chúng tôi cố gắng tìm hiểu các đặc trưng của các bộ dữ liệu tương ứng để cố gắng lý  giải thích  tại sao mô hình của chúng tôi hoặc các công trình khác có được kết quả tốt trên tập dữ liệu tương ứng đó.
% Những kết quả của hai đề xuất của chúng tôi cũng như các dịnh hướng nghiên cứu của chúng tôi trong tương lai.

Mặc dù kết quả chúng tôi cho thấy phương pháp dựa trên luật của chúng tôi có hiệu suất tương đương với các mô hình học sâu hiện đại (state-of-art) và có ưu thế vượt trội trong thời gian đào tạo khoảng 17 phút so với thời gian hàng giờ của phương pháp học sâu khác nhưng không phải là các mô hình học sâu này không đáng nghiên cứu. Chúng tôi cũng nhận thấy, với tập dữ liệu có nhiều loại quan hệ khác nhau như FreeBase, mô hình GCAT nhờ sử dụng cơ chế chú ý đạt được kết quả tốt hơn so với tập WorldNet với số lượng các loại quan hệ ít hơn.
Điều này chứng tỏ cơ chế chú ý bổ sung thêm thông tin vector nhúng quan hệ giúp học được các cấu trúc của đồ thị tốt hơn trên các tập dữ liệu nhiều loại quan hệ.
Đối với tập dữ liệu có nhiều mẫu tương tự và nghịch đảo như FB15k và WN18RR, mô hình dựa trên luật AnyBURL đạt được kết quả vượt trội, trong khi với phương pháp học sâu mô hình chỉ đạt kết quả trung bình so với các phương pháp khác.
Mô hình dựa trên luật AnyBURL có ưu thế tốt hơn trên các tập dữ liệu FB15k và WN18RR, tuy nhiên với các tập dữ liệu đã loại bỏ các thông tin tương tự hoặc nghịch đảo như FB15-237 và WN18RR thì phương pháp dựa trên học sâu đạt kết quả tốt hơn nhờ việc khái quát hóa được các cấu trúc phức tạp của đồ thị bằng việc cập nhật thông tin của một lượng lớn dữ liệu.
Một trong những ưu điểm của phương pháp dựa trên luật đó là các luật được sinh ra có thể lý giải được trong quá trình học, và có thời gian học vượt trội so với các phương pháp khác. Tuy nhiên, sau quá trình học phương pháp dựa trên luật phải duyệt qua tất cả các luật đã học mới có thể đưa ra được dự đoán. Đây là một điểm mà các phương pháp học sâu thể hiện tốt hơn, vì thông qua các trọng số đã học, mô hình học sâu GCAT thông qua các lớp tính toán có thể biến đổi đầu vào thành các kết quả xác xuất dự đoán nhanh hơn. Nhược điểm với phương pháp học sâu là quá trình học không thể lý giải được, hơn nữa chí phí huấn luyện rất tốn kém. Đối với hai thuật toán mở rộng của chúng tôi trong việc thêm tri thức mới vào đồ thị chúng tôi nhận thấy rằng là vượt trội hoàn toàn so với các phương pháp học sâu.
% Ngược lại đối với các phương pháp dựa trên học sâu lại có ưu thế rất lớn trong các tập dữ liệu này do có thể dễ dàng tính toán độ gần của các luật mới cần đánh giá so với các luật đã học từ đó có một kết quả khá tốt.
% Do đó chúng tôi cũng sẽ tiếp tục nghiên cứu các phương pháp học sâu và sẽ dùng phương pháp này làm đường cơ sở (base line) để so sánh với các nghiên cứu của chúng tôi trong tương lai.
% Một điểm yếu nữa của mô hình đựa trên luật của chúng tôi là mặc dù thời gian học là vượt trội nhưng thời gian để tính toán đưa ra đự đoán khá lâu do phải duyệt qua tất cả các luật được sinh ra mới có thể đưa ra dự đoán.
% Không giống như các phương pháp nhúng đồ thị khác thao tác này có thể dễ dàng tính toán.


Trong thế giới thực thông tin nhúng có thể thay đổi hoàn toàn theo thời gian, ví dụ trong quan hệ $\langle e_{\text{head}}, \textbf{president_of}, U.S\rangle$, năm 2012, ta biết kết quả là Barract Obama. Đối với nhiệm vụ dự đoán thực thể đầu, thì kết quả của $e_{\text{head}}$ có thể thay đổi hoàn toàn. Vì vậy, bổ sung thêm yếu tố thời gian là một hướng nghiên cứu quan trọng để đảm bảo tính đúng đắn của đồ thị tri thức. Ngoài ra, đối với phương pháp học sâu GCAT, có rất nhiều cải tiến trên cơ chế chú ý mà chúng tôi muốn áp dụng vào đồ thị tri thức. Đối với phương pháp học trên luật AnyBURL, gần đây nhánh học tăng cường (reinforcement learning) khá phát triển và nhóm tác giả Meilicke, Christian and Chekol \cite{meilicke2020reinforced} gần đây cũng đã có 1 nghiên cứu để tối ưu hóa lại phương pháp AnyBURL này.
Chúng tôi cũng có ý định nghiên cứu về hướng này và cố gắng báo cáo lại trong một tương lai gần.

% Công trình của tác giả (nếu không có thì comment 02 dòng dưới)
\addcontentsline{toc}{chapter}{Danh mục công trình của tác giả}
\chapter*{Danh mục công trình của tác giả}
\label{Publish}

\begin{enumerate}
\item Tạp chí ABC
\item Tạp chí XYZ
\end{enumerate}

% In tài liệu tham khảo
\addcontentsline{toc}{chapter}{Tài liệu tham khảo}
\printbibheading[title={Tài liệu tham khảo}]

\printbibliography[heading=subbibliography, title={Tiếng Việt}, keyword=Viet, resetnumbers=true]

\DeclareNameAlias{sortname}{last-first}
\DeclareNameAlias{default}{last-first}

\printbibliography[heading=subbibliography, title={Tiếng Anh}, notkeyword=Viet, resetnumbers=5] 
% ===================================================================== %
% CHÚ Ý: phải gán lại resetnumbers=số tài liệu tham khảo tiếng Việt + 1 %
% ===================================================================== %

% Phần phụ lục
\appendix 

\chapter{Các siêu tham số tối ưu}
\label{Appendix1}

Trong phần này chúng tôi sẽ báo cáo lại tập hợp các siêu tham số tối ưu trên cả mô hình Chú ý và mô hình ConvKB.
Chúng tôi sử dụng tìm kiếm lưới theo Hits@10 để tìm các tham số tối ưu. Chúng tôi không chia nhỏ tập dữ liệu trong quá trình huấn luyện mô hình chú ý,
đối với mô hình dự đoán, chúng tôi sử dụng kích thước cho tất cả các tập dữ liệu. Cụ thể các siêu tham số chúng tôi sử dụng được trình bày trong bảng sau :

\begin{table}[htbp]
	\begin{center}
		\resizebox{\textwidth}{!}{%
	\begin{tabular}{llllllllll}
		\hline
		& $\mu$  & Weight decay & Epochs & negative ratio  & Dropouts  & $\alpha_{\text{LeakyRELU}}$ & $N_{\text{head}}$ & $D_{\text{final}}$ & $\gamma$ \\
		\hline
		FB15k     & 1e-3 & 1e-5 & 3000 & 2 & 0.3 & 0.2 & 2 & 200 & 1 \\
		FB15k-237 & 1e-3 & 1e-5 & 3000 & 2 & 0.3 & 0.2 & 2 & 200 & 1 \\
		WN18      & 1e-3 & 5e-6 & 3600 & 2 & 0.3 & 0.2 & 2 & 200 & 5 \\
		WN18RR    & 1e-3 & 5e-6 & 3600 & 2 & 0.3 & 0.2 & 2 & 200 & 5 \\
		\hline
	\end{tabular}
}
\end{center}
\end{table}
\chapter{Bảng thuật ngữ}
\label{Appendix2}

\printglossary[type=\acronymtype]

\end{document} 